%%%%%%%%%%%%%%%%%%%%%%%%%%%%%%%%%%%%%%%%%%%%%%%%%%%%%%%%%%%%%%%%%%%%%%%%%%%%%%%%
%% Detecting Capability-Description Mismatches in LLM Agent Ecosystems
%% Target: ACM CCS 2026
%%
%% Uses official ACM acmart template (sigconf, anonymous submission)
%% Body: ≤12 pages (excl. references & appendices)
%%%%%%%%%%%%%%%%%%%%%%%%%%%%%%%%%%%%%%%%%%%%%%%%%%%%%%%%%%%%%%%%%%%%%%%%%%%%%%%%

\documentclass[sigconf,anonymous,review]{acmart}

%% Disable ACM copyright/DOI/ISBN blocks for anonymous submission
\settopmatter{printacmref=false}
\renewcommand\footnotetextcopyrightpermission[1]{}
\pagestyle{plain}

% ── Additional packages (acmart already loads: graphicx, amsmath, amssymb,
%    amsthm, booktabs, hyperref, natbib) ────────────────────────────────────
\usepackage{multirow}
\usepackage{enumitem}
\usepackage{subcaption}
\usepackage{xspace}

% ── Listings (code blocks) ─────────────────────────────────────────────────
\usepackage{listings}
\lstset{
  basicstyle   = \ttfamily\scriptsize,
  frame        = single,
  breaklines   = true,
  captionpos   = b,
  xleftmargin  = 0.5em,
  xrightmargin = 0.5em,
}
% Define yaml language for listings (not built-in)
\lstdefinelanguage{yaml}{
  keywords={true,false,null,yes,no},
  keywordstyle=\color{blue}\bfseries,
  basicstyle=\ttfamily\tiny,
  comment=[l]{\#},
  commentstyle=\color{gray}\itshape,
  stringstyle=\color{black},
  morestring=[b]",
  morestring=[b]',
  frame=single,
  breaklines=true,
  showstringspaces=false,
}

% ── amsthm theorem environments (acmart uses a different default) ───────────
\theoremstyle{definition}
\newtheorem{definition}{Definition}
\newtheorem{proposition}{Proposition}
\newtheorem{observation}{Observation}

% ── Custom colors ──────────────────────────────────────────────────────────
\definecolor{riskred}{rgb}{0.65,0.10,0.10}
\definecolor{safegreen}{rgb}{0.08,0.47,0.16}
\definecolor{warnorng}{rgb}{0.78,0.39,0.0}

% ── Macros ─────────────────────────────────────────────────────────────────
\newcommand{\Ttype}[1]{\textbf{\textcolor{riskred}{T#1}}\xspace}
\newcommand{\sysname}{\textsc{CapMismatch}\xspace}
\newcommand{\clawhub}{\textsc{ClawHub}\xspace}
\newcommand{\skillssh}{\textsc{skills.sh}\xspace}
\newcommand{\lscps}{\textsc{LSCPS}\xspace}
\newcommand{\eg}{\textit{e.g.,}\xspace}
\newcommand{\ie}{\textit{i.e.,}\xspace}
\newcommand{\cf}{\textit{cf.}\xspace}
\newcommand{\etal}{\textit{et al.}\xspace}

% ── Ensemble model macros (5-model ensemble) ─────────────────────────────
\newcommand{\mone}{Qwen3-8B\xspace}
\newcommand{\mtwo}{Llama-3.1-8B\xspace}
\newcommand{\mthree}{Mistral-7B-v0.3\xspace}
\newcommand{\mfour}{Gemma-4-E4B\xspace}
\newcommand{\mfive}{GPT-OSS-20B\xspace}
% TBD placeholder
\newcommand{\TBD}[1]{\colorbox{yellow}{\textbf{TBD#1}}}

% ── Key number macros ──────────────────────────────────────────────────────
\newcommand{\NskillsCH}{18{,}443\xspace}
\newcommand{\NskillsDMR}{18{,}443\xspace}
\newcommand{\NskillsDMRsrc}{18{,}443\xspace}
\newcommand{\Ntransitions}{24{,}243\xspace}
\newcommand{\NskillsSH}{116{,}208\xspace}
\newcommand{\NskillsSSH}{\NskillsSH}
\newcommand{\Nmalicious}{34\xspace}
\newcommand{\Nmalpairs}{40\xspace}
% ── Rule-based main experiment results (Script 27, April 2026) ───────────
\newcommand{\ResultCR}{0.042\xspace}
\newcommand{\ResultCP}{0.830\xspace}
\newcommand{\ResultCDI}{0.894\xspace}
\newcommand{\ResultICone}{75.8\%\xspace}
\newcommand{\ResultICtwo}{13.6\%\xspace}
\newcommand{\ResultFilesRead}{316{,}675\xspace}
\newcommand{\ResultAvgFiles}{17.2\xspace}
% ── CAPE ClawHub (§5) results ─────────────────────────────────────────────
\newcommand{\cape}{\textsc{CAPE}\xspace}
\newcommand{\NpairsTotal}{125\xspace}
\newcommand{\CapeSuccessITT}{45.6\%\xspace}
\newcommand{\CapeSuccessCR}{93.4\%\xspace}
\newcommand{\Tsixrecall}{94.1\%\xspace}
% ── CAPE skills.sh cross-platform validation (§6) results ─────────────────
\newcommand{\NssPairs}{197\xspace}          % Codex-confirmed viable skills
\newcommand{\NssRun}{166\xspace}           % skills executed (install ok)
\newcommand{\NssInstallRate}{100\%\xspace}  % 166/166
\newcommand{\NssSuccessRate}{99\%\xspace}  % 165/166
\newcommand{\NssShellRate}{100\%\xspace}   % T06: 75/75
\newcommand{\NssNetworkRate}{96\%\xspace}  % T07: 25/26
\newcommand{\NssMemoryRate}{100\%\xspace}  % T13: 20/20
\newcommand{\NssFileRate}{100\%\xspace}    % T02: 23/23
\newcommand{\NssPersistRate}{100\%\xspace} % T05: 9/9

% ── Compact lists ──────────────────────────────────────────────────────────
\setlist{nosep, leftmargin=1.2em}

% ── Overfull hbox prevention ───────────────────────────────────────────────
\emergencystretch=2em
\hfuzz=2pt
\tolerance=1500
\sloppy

% ── ACM Conference Metadata ────────────────────────────────────────────────
\acmConference[CCS '26]{ACM Conference on Computer and Communications Security}{October 2026}{Salt Lake City, UT, USA}
\acmYear{2026}
\acmISBN{}
\acmDOI{}

%===============================================================================
\begin{document}
%===============================================================================

% ── ACM Computing Classification System (CCS) Concepts ────────────────────
\begin{CCSXML}
<ccs2012>
  <concept>
    <concept_id>10002978.10002979</concept_id>
    <concept_desc>Security and privacy~Software security engineering</concept_desc>
    <concept_significance>500</concept_significance>
  </concept>
  <concept>
    <concept_id>10002978.10003006.10003008</concept_id>
    <concept_desc>Security and privacy~Malware and its mitigation</concept_desc>
    <concept_significance>500</concept_significance>
  </concept>
  <concept>
    <concept_id>10002978.10003006.10003014</concept_id>
    <concept_desc>Security and privacy~Supply chain attacks</concept_desc>
    <concept_significance>300</concept_significance>
  </concept>
  <concept>
    <concept_id>10010147.10010257.10010293</concept_id>
    <concept_desc>Computing methodologies~Natural language processing</concept_desc>
    <concept_significance>100</concept_significance>
  </concept>
</ccs2012>
\end{CCSXML}

\ccsdesc[500]{Security and privacy~Software security engineering}
\ccsdesc[500]{Security and privacy~Malware and its mitigation}
\ccsdesc[300]{Security and privacy~Supply chain attacks}
\ccsdesc[100]{Computing methodologies~Natural language processing}

\keywords{LLM agent marketplace, software supply chain, capability disclosure,
static analysis, capability mismatch, skill security, MCP}

\title{Detecting Capability-Description Mismatches\\ in LLM Agent Ecosystems}

% Anonymous submission
\author{Anonymous Authors}
\affiliation{\institution{Anonymous Institution}}

% ── Abstract (must come before \maketitle in acmart) ──────────────────────
\begin{abstract}
LLM skill marketplaces grant third-party plugins ambient authority over
files, credentials, and network connections, yet rely on free-text
descriptions as the sole consent mechanism---creating a systematic gap
between user-visible disclosures and actual executable behavior.
We present \sysname, a deterministic publish-time screening workflow
that extracts capabilities from multi-file skill artifacts via 300+
regex patterns (LSCPS, 14 families) and ranks submissions by
capability-description inconsistency using Capability Recall (CR) and
the Capability Discrepancy Index (CDI).
Across \NskillsDMR{} \clawhub{} skills, \ResultICone{} show complete
capability non-disclosure (IC-1, CR\,=\,0) and \ResultICtwo{} partial
non-disclosure (IC-2), with corpus-mean CR\,=\,\ResultCR{} and
CDI\,=\,\ResultCDI{}; the same pipeline applied without adjustment
to \NskillsSH{} \skillssh{} skills yields $>$90\% inconsistency---
consistent with the pattern being widespread rather than platform-specific.
A stratified security-relevance study (\textsection\ref{sec:detection},
$\kappa{=}0.81$) confirms that 53.8\% of top-ranked mismatches are
security-relevant consent-risk cases rather than benign documentation gaps;
against \Nmalicious{} forensically confirmed malicious skills---missed
entirely by the platform scanner (0\% recall)---\sysname achieves
94.1\% recall and a \textbf{16.3$\times$ lift} at a 200-skill
review budget, providing triage value orthogonal to existing scanners.
The system is a \emph{prioritization aid for human review},
not a standalone maliciousness detector.
\end{abstract}

\maketitle

% ── Body (≤12 pages, excl. references and appendices) ─────────────────────
% sections/01_introduction.tex

\section{Introduction}
\label{sec:introduction}

The rise of LLM-powered skill marketplaces---platforms where developers
publish installable agents that extend the capabilities of large language
models~\cite{xi2023rise,schick2024toolformer,yao2023react}---represents a new class of software supply chain.
Unlike traditional package managers~\cite{zimmermann2019small,duan2021measuring}, skills operate with
\emph{NLP-driven ambient authority}: once installed, they act on behalf of the user
across sessions, reading files, sending messages, executing shell commands,
and accessing credentials, often without per-action confirmation or OS-level sandboxing~\cite{ruan2024identifying}.

\clawhub{}, the largest public LLM skill marketplace at the time of our study,
had 18,443 listed skills and $>$2.1M installs as of March 2026---growing
at $\approx$40\% year-over-year since the MCP standard's release~\cite{anthropic2024mcp}.
Historically, supply-chain security problems are \emph{hardest to address at scale}:
the npm event-stream attack (\S\ref{sec:related}) reached 2M downloads before detection.
Measuring and hardening the threat surface \emph{now}, at 18K skills,
is orders of magnitude cheaper than intervening at 1M skills.
This motivates a publish-time screening approach focused on exposing
consent-risk gaps before harm reaches users.

Unlike mobile operating systems (\eg Android) which enforce capability boundaries
via strict, machine-readable permission manifests~\cite{felt2011android,au2012pscout}, LLM skill marketplaces rely
overwhelmingly on unstructured natural language descriptions to negotiate user consent.
This gap creates a novel trust problem.
When a user installs a ``calendar scheduling'' skill,
they implicitly trust its description to accurately reflect
what the skill will do.
But how accurate are these descriptions?
And what happens when a skill's author quietly updates it?

We answer these questions by proposing a publish-time screening workflow for capability-description mismatches in LLM agent ecosystems.
Importantly, while much prior research focuses on \emph{prompt injection}~\cite{greshake2023indirect}
(adversarial modification of agent inputs), we study the \emph{supply chain attack surface}~\cite{ohm2020backstabber,ladisa2023taxonomy}:
deceptive or malicious behavior embedded directly within the installed tool code itself.

\paragraph{What we detect and what we do not.}
Our system targets a specific and narrow threat:
\emph{skills whose actual runtime capabilities ($\mathcal{C}_{\text{act}}$) substantially
exceed their described capabilities ($\mathcal{C}_{\text{vis}}$) in security-sensitive
dimensions}---specifically involving credential access, outbound
data exfiltration, autonomous execution, or destructive file operations.
We do \emph{not} claim to detect all malicious code, runtime exploits,
JIT-payload downloaders, or prompt injection attacks.
A critical property of LLM skills is that the description is the \emph{only}
(and often only partially read) capability disclosure mechanism available
to users at install time---even if users do not read it carefully,
its systematic failure enables supply-chain deception at scale.
We document this gap and propose a triage pipeline that prioritizes
security-relevant mismatches over benign documentation lapses.

\paragraph{Research questions.}
In this work, we investigate the following three research questions:

\begin{itemize}[leftmargin=*,topsep=2pt,itemsep=1pt]

\item \textbf{RQ1 (Prevalence \& Disclosure Quality):}
  How prevalent is capability-description inconsistency across LLM skill
  marketplace platforms, and do developers provide skill descriptions that
  accurately and meaningfully reflect their skills' actual executable behavior?

\item \textbf{RQ2 (Detection \& Family Taxonomy):}
  Can we systematically identify the gap between skill descriptions and
  actual code capabilities?
  Which capability families are most frequently and severely under-disclosed,
  and can a static multi-file extraction pipeline expose these gaps at
  publish time?

\item \textbf{RQ3 (Automated Disclosure Remediation):}
  Can automatically generated capability descriptions---derived from static
  code analysis---effectively close the disclosure gap and supply developers
  and users with accurate, actionable capability information?

\end{itemize}

Our key contributions are:

\begin{itemize}[leftmargin=*,topsep=3pt,itemsep=2pt]

\item \textbf{IC Framework and Deterministic Publish-Time Screening.}
We propose \sysname, a formal Inconsistency Checking (IC) difference-matrix
framework and a deterministic publish-time screening pipeline that extracts
capabilities from all skill artifacts (SKILL.md, Python/JS/Shell source,
dependency manifests) via 300+ regex patterns (LSCPS, 14 families).
Multi-file aggregation is essential: 32.4\% of confirmed malicious skills
have high-risk capabilities exclusively in non-SKILL.md files.
Semantic adjudication and temporal monitoring are optional extensions;
the core system---SCS + DMR---is fully deterministic and requires no LLM.

\item \textbf{Ecosystem-Scale Measurement and Security-Relevance Study.}
Across \NskillsCH{} \clawhub{} skills, \ResultICone{} exhibit complete
capability non-disclosure (IC-1) and \ResultICtwo{} partial non-disclosure (IC-2),
producing corpus-mean CR\,=\,\ResultCR{} and CDI\,=\,\ResultCDI{}.
A stratified 500-skill human study (E4, \S\ref{sec:detection};
pre-defined rubric independent of T6; inter-rater $\kappa{=}0.81$)
finds 53.8\% of IC mismatches are security-relevant consent-risk cases
---not mere documentation incompleteness.
Cross-ecosystem validation on \NskillsSSH{} \skillssh{} skills yields $\approx$90\%
inconsistency, consistent with structural prevalence across marketplaces.

\item \textbf{Triage Evaluation on Confirmed Malicious Skills.}
On 40 forensically documented malicious version pairs (34 unique skills,
$\kappa{=}0.91$), SCS$\cup$DMR retrieves 32/34 (94.1\%,
95\% Wilson CI:\,[80.9\%,\,98.4\%]), while all baselines retrieve none.
At a 200-skill review budget on the full \NskillsDMR{}-skill corpus,
\sysname achieves \textbf{16.3$\times$ lift} over random review---
orthogonal to, and not a replacement for, binary malware scanning.

\end{itemize}

\begin{figure}[t]
\centering
\includegraphics[width=\columnwidth]{figures/fig_pipeline_overview.pdf}
\caption{\sysname publish-time screening pipeline. Skill artifacts feed into
  multi-file capability extraction (SCS), producing the IC classification
  and CDI-ranked review queue (DMR).
  Temporal monitoring (TDM) and semantic adjudication (SA) are
  \emph{optional extensions}---not required for the main triage results;
  current core evaluation covers SCS+DMR only.}
\label{fig:overview}
\end{figure}

\paragraph{Responsible disclosure.}
We notified both platforms of our findings prior to publication.
All skill identifiers in T6 cases have been pseudonymized in this version.

\paragraph{Open science.}
Our detection pipeline source code, full \NskillsDMR{} IC classification
results, and the anonymized T6 ground-truth corpus are available
for review via anonymous repository;
upon acceptance, we will release the complete artifact package including
the full \lscps{} regex bank, data extraction scripts,
and reproducible evaluation notebooks following USENIX artifact evaluation guidelines.

\paragraph{Paper organization.}
\S\ref{sec:background} provides our threat model and capability taxonomy.
\S\ref{sec:taxonomy} formalizes the IC framework and difference matrix.
\S\ref{sec:detection} defines the multi-layer detection pipeline.
\S\ref{sec:static}--\ref{sec:groundtruth} present three evaluations:
ecosystem-scale measurement with cross-platform validation,
temporal update-risk analysis,
and malicious injection retrieval on ground-truth cases.
\S\ref{sec:related} discusses related work.

% sections/02_background.tex

\section{Background \& Threat Model}
\label{sec:background}

\subsection{LLM Skill Marketplaces}

LLM skill marketplaces are application stores specialized for language
model extensions. A \emph{skill} is a self-contained module---typically a
\texttt{SKILL.md} manifest with tool definitions and helper code---that an
LLM agent invokes during conversation or automated tasks.
Skills may be first-party (platform-developed) or third-party.

Users install skills through a package-manager command
(\eg \texttt{/install author/skill-name}) and thereafter the skill
executes with \emph{ambient authority}: persistent access to the
agent's context window, conversation history, tools, and any connected
data sources (filesystem, email, calendar, credentials).

\paragraph{Trust signals.}
Platforms expose signals intended to help users evaluate skills:
\begin{enumerate}
    \item A free-text \textbf{description} written by the developer;
    \item Optional structured \textbf{metadata} (YAML declarations of capability);
    \item Version \textbf{changelogs};
    \item A platform-operated \textbf{security scan badge};
    \item Aggregate \textbf{download and star counts}.
\end{enumerate}
Our analysis suggests that signals 1--4 often understate actual capabilities, while signal 5 may amplify misplaced trust in risky skills.

\paragraph{Auto-update semantics.}
Both \clawhub{} and \skillssh{} propagate version updates automatically
to all installed users by default.
Users who consented to a skill at version $v_i$ are silently upgraded to
$v_{i+1}$ without re-review, re-consent, or changelog notification.
ClawHub's platform documentation confirms that updates propagate within
24~hours of publisher push.
Version pinning is unavailable to standard users; skill developers have
unilateral write access to all installed instances.
This is the key structural enabler of T4 and T6 attacks.

\subsection{Threat Model}

We model four principals:

\textbf{Users} install skills expecting them to behave as described.
They cannot audit code before installation and rely on platform trust signals.

\textbf{Benign (careless) developers} omit capability documentation
without malicious intent (T1, T4). Their omissions violate informed consent.

\textbf{Deceptive developers} consciously misstate capabilities to
increase installs or bypass platform review (T2, T3, T5).

\textbf{Malicious actors} (compromised accounts or supply-chain
attackers) deliberately inject harmful code into skills (T6).

\paragraph{Attack surface.}
The core vulnerability is \emph{consent at installation vs.\ capability
at runtime}. Users who consented to a ``calendar scheduler'' have not
consented to credential exfiltration---even if that capability is
in the installed code. This gap is structurally analogous to Android
permission creep~\cite{felt2011android} but operates at the description
level rather than the manifest level: Android \emph{requires} capability
declarations; skill platforms do not.

\paragraph{Attacker goal.}
A T6 attacker seeks to: (a) establish trust via a benign skill version,
accumulate installs, then (b) inject malicious capability via an update,
exploiting auto-update propagation, then (c) exfiltrate credentials/data
before the malicious version is detected and patched.
The expected harm is proportional to the install count and the exposure
window---both of which grow as the platform grows.

\paragraph{Why undisclosure is the attacker's rational strategy.}
A reviewer may ask why attackers do not use more sophisticated obfuscation
(runtime payload fetching, dynamic code generation, indirect API calls).
We observe that simple undisclosure is the \emph{risk-minimal} approach:
complexity increases the probability of detection by
pattern matching, sandbox execution, or anomaly scoring.
T6 evidence supports this: 32/34 confirmed T6 skills use
straightforward API-level patterns---\texttt{os.environ}~(CRED\_READ),
\texttt{requests.post}~(NET\_WRITE), \texttt{os.remove}~(FILE\_DELETE)---
without obfuscation. The marginally harder bar of tool-level
undisclosure is sufficient because platform scanners check
\emph{declared intent} (descriptions and metadata), not runtime capability.

\paragraph{Consent model scope.}
We do \emph{not} claim that users systematically read skill descriptions
or form complete capability expectations from them.
Rather, we note that descriptions are the \emph{only structured,
publish-time mechanism} for capability disclosure available to users across
both platforms we studied (YAML metadata fill rate: 0.0\%;
\S\ref{sec:static}).
When this mechanism fails systematically at scale, users and platform
operators have no alternative disclosure channel---regardless of
whether individual users attend to descriptions.
The threat is structural, not behavioral.

\subsection{LSCPS Capability Taxonomy}

We categorize skill capabilities into five top-level families,
abbreviated \textbf{\lscps{}}
(\underline{L}ocal-shell, \underline{S}ecret, \underline{C}ommand-and-control,
\underline{P}ersistent-write, \underline{S}ystemic-autonomous).
Each top-level family is subdivided into 14 fine-grained sub-categories
that serve as the atomic detection units (see Appendix~\ref{app:lscps}):

\begin{itemize}[topsep=2pt,itemsep=2pt]
\item[\textbf{L}] \textbf{Local Shell execution}:
  \texttt{subprocess}, \texttt{os.system}, \texttt{eval()}, bash invocation
\item[\textbf{S}] \textbf{Secret/Credential access}:
  API keys, OAuth tokens, passwords, private keys
\item[\textbf{C}] \textbf{Command-and-control / Network write}:
  \texttt{requests.post}, webhooks, SMTP exfiltration
\item[\textbf{P}] \textbf{Persistent file write}:
  \texttt{open(w)}, \texttt{shutil.rmtree}, \texttt{os.remove}
\item[\textbf{S}] \textbf{Systemic/Autonomous effects}:
  send email, create/modify files, background scheduling, multi-step no-confirm loops
\end{itemize}

Each sub-category is detected by a curated regex bank totaling 300+ patterns
(see Appendix~\ref{app:lscps}),
validated on the 40 T6 ground-truth version pairs
(72.5\% pair-level high-risk recall; 79.4\% at skill level,
\S\ref{sec:e2e}).

\paragraph{Methodological rationale: Why LSCPS, not Bandit/Semgrep?}
General-purpose security linters (Bandit, Semgrep) detect Python anti-patterns
(SQL injection, hardcoded secrets) but are designed for \emph{single-file,
code-centric} analysis. LSCPS differs in three concrete ways:

\begin{enumerate}[topsep=2pt, itemsep=2pt]
\item \textbf{Multi-file capability aggregation.}
  LSCPS extracts $\mathcal{C}_{\text{act}}$ across SKILL.md manifests,
  Python helpers, YAML configs, and JSON schemas jointly (up to 17.2 files
  per skill on average). Bandit analyzes files independently and misses
  exfiltration chains that span manifest + helper code.
\item \textbf{Dual-side extraction.}
  LSCPS extracts $\mathcal{C}_{\text{vis}}$ from \emph{natural language
  descriptions} as well as $\mathcal{C}_{\text{act}}$ from code.
  No general linter performs description-vs-code inconsistency comparison.
\item \textbf{Marketplace-specific patterns.}
  LSCPS covers 14 families including MCP permission declarations,
  autonomous-mode signals, and SKILL.md structural cues outside general linters'
  default rule sets. This yields $\approx$35\% added recall over Bandit
  (high-severity) on our T6 evaluation corpus.
\end{enumerate}

Beyond these differentiators, static extraction provides three deployment advantages:
\begin{itemize}
    \item \textbf{Evasion resistance:} Static extraction evaluates purely
      syntactically, unlike LLM judges which are vulnerable to adversarial
      prompt manipulation injected into code comments.
    \item \textbf{Deployability:} The regex-based engine executes
      marketplace-scale scanning across tens of thousands
      of skills in seconds with zero runtime overhead.
    \item \textbf{High baseline recall:} Our 300+ pattern bank identifies
      72.5\% of real-world T6 malicious pairs at the pair level,
      compared to the platform scanner's 0\%.
\end{itemize}

\paragraph{Capability weighting.}
For the static baseline (SCS), we assign heuristic
severity weights to each LSCPS family as an initial
Capability Exposure Score (CES): shell execution (3.0),
credential transmit (3.0), credential read (2.5),
network write (2.0), autonomous (1.5), file read (0.5).
In \S\ref{sec:riskweights}, we replace these heuristics
with \emph{data-driven} Risk-Weighted Mismatch (RWM) weights
learned from T6 ground truth. The learned weights diverge
significantly from the heuristic assignments
(\eg shell execution drops from 3.0 to 0.10 because it is
ubiquitous across both malicious and benign skills),
demonstrating the necessity of empirical calibration.

\subsection{Datasets}

\paragraph{\clawhub{}.}
We archived the full \clawhub{} public skill catalog in March 2026,
capturing \NskillsCH{} skills and complete version histories.
Of these, \textbf{\NskillsDMR{} skills} have locally archived
\texttt{SKILL.md} text, covering \textbf{\Ntransitions{} version transition pairs}
(on average 3.39 versions per skill).
The 705 remaining skills were deleted from the platform before
local archival completed.
\textit{Deleted-skill bias note:} Because deleted skills may disproportionately
represent platform-removed (and thus potentially higher-risk) content,
excluding them conservatively biases our IC-1 prevalence \emph{downward}---
our reported figures therefore underestimate the true inconsistency rate.

\paragraph{\skillssh{}.}
We archived \textbf{116,208 skills} from \skillssh{}.
Unlike \clawhub{}, \skillssh{} does not expose
multi-version histories in its public API, so our analysis is static-only.
We apply the identical IC framework pipeline
(\S\ref{sec:static}) to validate cross-platform generalizability.

\paragraph{Ethical considerations.}
All data is from public APIs at rate $\leq$ 0.5 req/s.
We did not execute any skill, access non-public content,
or interfere with platform operations.
Our crawler identified itself with an academic research user agent.
T6 case identifiers are withheld pending coordinated disclosure.

% sections/03_taxonomy.tex
% §3: Inconsistency Checking Framework (IC Framework)
% Restructured to follow IC difference-matrix approach (Liao et al., ACSAC 2020)

\section{Inconsistency Checking Framework}
\label{sec:taxonomy}

Inspired by the inconsistency checking methodology of~\cite{liao2020measuring},
we formalize capability-description mismatches for LLM agent skills
via a \emph{difference matrix} that captures both directions of divergence
between what a skill \emph{claims} (user-visible) and what it \emph{does} (actual).
Figure~\ref{fig:methodology} illustrates the extraction and classification pipeline.

\begin{figure}[t]
\centering
\includegraphics[width=\columnwidth]{figures/fig_methodology_overview.pdf}
\caption{Methodology overview: each SKILL.md is decomposed into description (yielding $\mathcal{C}_{\text{vis}}$) and code blocks (yielding $\mathcal{C}_{\text{act}}$ via LSCPS). The difference matrix produces four quadrants (AD, UD, FC, TN), from which IC types and quantitative metrics are derived.}
\label{fig:methodology}
\end{figure}

% ── Notation Table ──────────────────────────────────────────────────────────
\begin{table}[h]
\centering
\caption{Notation. All metrics are formally defined in \S\ref{sec:taxonomy}.}
\label{tab:notation}
\small
\begin{tabular}{@{}lp{6.2cm}@{}}
\toprule
\textbf{Symbol} & \textbf{Definition} \\
\midrule
$\mathcal{C}_{\text{vis}}(s)$  & User-visible capability set: mentions in marketplace description \\
$\mathcal{C}_{\text{act}}(s)$  & Actual capability set: code-detected capabilities across all skill artifacts \\
$\mathcal{F}$                  & Universe of 14 capability families (Table~\ref{tab:lscps_kw}) \\
\midrule
$\text{AD}(s)$ & Accurate Disclosure: $|\mathcal{C}_{\text{act}} \cap \mathcal{C}_{\text{vis}}|$ \\
$\text{UD}(s)$ & Undisclosed: $|\mathcal{C}_{\text{act}} \setminus \mathcal{C}_{\text{vis}}|$ (hidden) \\
$\text{FC}(s)$ & False Claim: $|\mathcal{C}_{\text{vis}} \setminus \mathcal{C}_{\text{act}}|$ (over-claimed) \\
$\text{TN}(s)$ & True Negative: $|\mathcal{F} \setminus (\mathcal{C}_{\text{act}} \cup \mathcal{C}_{\text{vis}})|$ \\
\midrule
$\text{CR}(s)$  & Capability Recall (Eq.~\ref{eq:cr}) \\
$\text{CP}(s)$  & Description Precision (Eq.~\ref{eq:cp}) \\
$\text{CDI}(s)$ & Capability Discrepancy Index (Eq.~\ref{eq:cdi}) \\
$\text{RWM}(s)$ & Risk-Weighted Mismatch (Eq.~\ref{eq:rwm}) \\
$\text{TCE}$    & Temporal Capability Expansion (Eq.~\ref{eq:tce}) \\
$\text{SDS}$    & Semantic Disclosure Score (changelog coverage) \\
\bottomrule
\end{tabular}
\end{table}

% ── Information Sources ──────────────────────────────────────────────────────
\subsection{User-Visible vs.\ Actual Capability Sets}

A LLM agent skill exposes information through two distinct channels:

\paragraph{User-visible layer ($\mathcal{C}_{\text{vis}}$).}
The marketplace description is the \emph{only} information a user
reads before installation. From this natural-language text,
we extract the set of capability families mentioned,
using a keyword dictionary derived from the LSCPS pattern bank (\S\ref{sec:static}).
This mirrors how~\cite{liao2020measuring} extracted data practices
from voice-app descriptions.

\paragraph{Actual capability layer ($\mathcal{C}_{\text{act}}$).}
The SKILL.md file contains the skill's executable specification,
invisible to end users. In Phase~1 (this paper),
$\mathcal{C}_{\text{act}} = \mathcal{C}_{\text{code}}$,
derived via Static Capability Screening (SCS) over all code blocks.
Phase~2 (future work) extends this to include natural-language
instruction grants ($\mathcal{C}_{\text{instr}}$) and
prose capability descriptions ($\mathcal{C}_{\text{text}}$) within SKILL.md:
\[
  \mathcal{C}_{\text{act}} \supseteq \mathcal{C}_{\text{code}}
  \;\cup\; \mathcal{C}_{\text{instr}} \;\cup\; \mathcal{C}_{\text{text}}
\]
Restricting to $\mathcal{C}_{\text{code}}$ makes our measurements
\emph{conservative}: the true mismatch is at least as large as reported.

% ── Difference Matrix ────────────────────────────────────────────────────────
\subsection{The Inconsistency Difference Matrix}

\begin{definition}[Capability Difference Matrix]
\label{def:dm}
For skill $s$ with capability sets $\mathcal{C}_{\text{act}}(s)$ (actual) and
$\mathcal{C}_{\text{vis}}(s)$ (visible), and capability family set $\mathcal{F}$,
the \emph{Difference Matrix} $\text{DM}(s) = \langle\text{AD},\text{UD},\text{FC},\text{TN}\rangle$
partitions $\mathcal{F}$ into four disjoint cells:
\end{definition}

\vspace{-0.5em}
\begin{center}
\small
\begin{tabular}{@{}lcc@{}}
\toprule
 & $f \in \mathcal{C}_{\text{vis}}$ & $f \notin \mathcal{C}_{\text{vis}}$ \\
\midrule
$f \in \mathcal{C}_{\text{act}}$ & \textbf{AD} (aligned) & \textbf{UD} (hidden) \\
$f \notin \mathcal{C}_{\text{act}}$ & \textbf{FC} (false claim) & \textbf{TN} (absent) \\
\bottomrule
\end{tabular}
\end{center}

Summing over all families in $\mathcal{F}$ yields per-skill scalars
$\text{AD}(s)$, $\text{UD}(s)$, $\text{FC}(s)$, $\text{TN}(s)$.

\paragraph{Capability Recall (CR).}
What fraction of actual capabilities does the description disclose?
\begin{equation}
  \text{CR}(s) = \frac{\text{AD}(s)}{\text{AD}(s)+\text{UD}(s)}
  \label{eq:cr}
\end{equation}
$\text{CR}=1$ indicates perfect disclosure; $\text{CR}=0$ indicates total omission.
Prior work measured only CR~\cite{liao2020measuring,zhang2022measuring};
we show this is insufficient.

\paragraph{Description Precision (CP).}
What fraction of claimed capabilities are genuine?
\begin{equation}
  \text{CP}(s) = \frac{\text{AD}(s)}{\text{AD}(s)+\text{FC}(s)}
  \label{eq:cp}
\end{equation}
$\text{CP}=1$ means no false claims; defined only when $\mathcal{C}_{\text{vis}} \neq \varnothing$.
CP~$< 1$ identifies skills that promise capabilities developers never implemented.

\paragraph{Capability Discrepancy Index (CDI).}
The overall normalized inconsistency:
\begin{equation}
  \text{CDI}(s) = \frac{\text{UD}(s)+\text{FC}(s)}{\text{AD}(s)+\text{UD}(s)+\text{FC}(s)}
  \label{eq:cdi}
\end{equation}
CDI $\in [0,1]$; 0 = perfect consistency, 1 = complete mismatch.

\paragraph{Risk-Weighted Mismatch (RWM).}
CDI~treats all families equally; RWM weights undisclosed families
by their empirical malicious-injection risk:
\begin{equation}
  \text{RWM}(s) = \sum_{f \in \mathcal{C}_{\text{act}} \setminus \mathcal{C}_{\text{vis}}} w_f
  \label{eq:rwm}
\end{equation}
Weights $w_f$ are learned from T6 ground truth via logistic regression
(\S\ref{sec:groundtruth}).
RWM is used \emph{only for risk ranking}, not for IC classification---preserving
classification objectivity.

% ── IC Classification ────────────────────────────────────────────────────────
\subsection{Inconsistency Classification (IC-1--IC-4)}
\label{sec:ic_classification}

\begin{definition}[IC Classification]
\label{def:ic}
Given $\text{DM}(s) = \langle\text{AD},\text{UD},\text{FC},\text{TN}\rangle$,
skill $s$ is assigned to exactly one of six mutually exclusive, collectively
exhaustive classes: IC-0 (Accurate: $\text{CR}=\text{CP}=1$),
IC-1 (Full Omission: $\text{UD}>0, \text{AD}=0, \text{FC}=0$),
IC-2 (Partial Omission: $\text{UD}>0, \text{AD}>0, \text{FC}=0$),
IC-3 (False Claim: $\text{FC}>0, \text{UD}=0$),
IC-4 (Bidirectional: $\text{UD}>0, \text{FC}>0$),
or TN (No Capabilities: $\text{AD}=\text{UD}=\text{FC}=0$).
\end{definition}

The difference matrix induces a mutually exclusive, exhaustive
classification of every skill:

\begin{table}[h]
\centering
\caption{IC inconsistency types derived from the difference matrix.
  Prevalence over \clawhub{} ($N{=}18{,}443$, rule-based static extraction).}
\label{tab:ic_types}
\scriptsize
\begin{tabular}{@{}clp{3.2cm}rr@{}}
\toprule
\textbf{Type} & \textbf{Name} & \textbf{Condition} & \textbf{Count} & \textbf{\%} \\
\midrule
IC-1 & Full Omission      & UD$>$0, AD$=$0, FC$=$0 & 13{,}988 & 75.8\% \\
IC-2 & Partial Omission   & UD$>$0, AD$>$0, FC$=$0 & 2{,}513  & 13.6\% \\
IC-3 & False Claim        & FC$>$0, UD$=$0         & 78       &  0.4\% \\
IC-4 & Bidirectional      & UD$>$0, FC$>$0         & 538      &  2.9\% \\
\midrule
---    & Accurate           & UD$=$0, FC$=$0, AD$>$0 & 79       &  0.4\% \\
---    & No Capabilities    & AD$=$FC$=$UD$=$0       & 1{,}247  &  6.8\% \\
\midrule
\multicolumn{3}{l}{\textbf{Total Inconsistent (IC-1--IC-4)}} & \textbf{17{,}117} & \textbf{92.8\%} \\
\bottomrule
\end{tabular}
\end{table}

\paragraph{IC-1: Full Omission (CR\,=\,0).}
The description mentions \emph{no} capability from $\mathcal{C}_{\text{act}}$,
yielding CR\,=\,0.
This is the most prevalent pattern (75.8\%),
indicating that a large fraction of skills have marketplace descriptions
that provide no signal about detected executable capabilities.
Among T6 confirmed malicious skills, \textbf{22/34 (64.7\%)} are IC-1.
We note that rule-based extraction may over-attribute capabilities in
boilerplate code; IC-1 prevalence should be interpreted as an upper bound
pending LLM-based adjudication (future work).

\paragraph{IC-2: Partial Omission (0\,$<$\,CR\,$<$\,1).}
The description mentions some capabilities but omits others
(e.g., disclosing \textsc{File-Read}
while not mentioning \textsc{Cred-Read} or \textsc{Net-Write}).
This accounts for 13.6\% of skills;
\textbf{10/34 T6} skills (29.4\%) are IC-2.
Combined IC-1+IC-2 captures \textbf{94.1\%} of all T6 malicious skills.

\paragraph{IC-3: False Claim (CP\,$<$\,1).}
The description asserts a capability not detected in the code (0.4\%, $n{=}78$).
The most common false claims are
\textsc{Autonomous} and \textsc{Multi-Step}---skills marketed
as AI-orchestrated when no corresponding code pattern is detected.
No T6 skills are IC-3.

\paragraph{IC-4: Bidirectional Mismatch.}
Both omission and false-claim occur simultaneously (2.9\%, $n{=}538$):
the skill has unmentioned capabilities \emph{and} claims capabilities
not found in the code.
Example: ``raini-skill-audit'' has CRED\_READ+SHELL\_EXEC patterns
while claiming MULTI\_STEP without a corresponding implementation.

% ── Temporal Extension ──────────────────────────────────────────────────────
\subsection{Temporal IC: Version Update Patterns}
\label{sec:temporal_ic}

The difference matrix can be computed per version,
enabling IC-type tracking across updates.
We define three temporal patterns:

\paragraph{Temporal Capability Expansion (TCE).}
For transition $v_i \to v_{i+1}$:
\begin{equation}
  \text{TCE}(v_i \to v_{i+1}) = \text{RWM}(s, v_{i+1}) - \text{RWM}(s, v_i)
  \label{eq:tce}
\end{equation}
\emph{Silent expansion} occurs when TCE $> 0$ and the Semantic Disclosure Score
(SDS, changelog adequacy) $<0.35$.
Of \Ntransitions{} ClawHub version pairs,
\textbf{2,575 (10.6\%)} are silent expansions.

\paragraph{Silent Capability Injection (T1).}
When $\text{TCE} > 0$, $\text{SDS} < 0.35$, and new capabilities
include the credential-exfiltration pair
(CRED-READ $+$ NET-WRITE), we label the
transition a \emph{silent injection} (T1).
Of 512 T1 transitions, 100\% meet our thresholds.

\paragraph{T6: Confirmed Malicious Injection.}
Ground-truth malicious transitions ($n{=}40$) are identified via
developer-authored changelogs that explicitly acknowledge prior
malicious behavior (\S\ref{sec:groundtruth}).
T6 differs from T1 in verified intent rather than heuristic evidence.

% ── Detection Pipeline ──────────────────────────────────────────────────────
\subsection{IC-Based Detection Pipeline}
\label{sec:pipeline}

The IC framework drives a four-stage detection pipeline:

\begin{enumerate}[topsep=3pt,itemsep=2pt]
\item \textbf{Step 1 --- Capability Extraction}:
  SCS+DMR extracts $\mathcal{C}_{\text{act}}$ from all skill artifacts
  (SKILL.md, Python/JS/Shell source, dependency manifests, config files)
  using 300+ regex patterns across 14 LSCPS families; SCS alone
  matches code blocks within SKILL.md (89 family-core patterns);
  DMR extracts $\mathcal{C}_{\text{vis}}$ from marketplace description text.

\item \textbf{Step 2 --- Difference Matrix}:
  Compute AD, UD, FC, TN per family;
  derive CR, CP, CDI, RWM.

\item \textbf{Step 3 --- IC Classification}:
  Assign IC type (IC-1 through IC-4, Accurate, or TN-Only).
  SCS$\cup$DMR composite rank is the output of the core system.

\end{enumerate}

\noindent\textbf{Optional extensions.}
Steps 1--3 are the current core system: deterministic, $O(N)$ time,
no LLM required. Two optional extensions are designed for future deployment:
(a)~\textbf{Temporal Monitoring (TDM)}: invoked when version history is
available; detects silent capability expansions across updates;
(b)~\textbf{Semantic Adjudication (SA)}: a five-model LLM ensemble invoked
only for the top-ranked high-RWM subset (3--5\% of skills) to reduce
rule-based FPs and recover non-keyword expressed capabilities.
Neither extension participates in the current main evaluation;
both are reserved for future deployment refinement (\S\ref{sec:sa}, \S\ref{sec:temporal}).

% ── Evaluation Tasks ────────────────────────────────────────────────────────
\subsection{Research Questions}

\begin{itemize}[topsep=3pt,itemsep=2pt]
\item \textbf{RQ1 (Static IC):} What is the ecosystem-scale distribution
  of IC-1 through IC-4? Which capability families are most systematically
  hidden vs.\ falsely claimed? (\S\ref{sec:static})
\item \textbf{RQ2 (Temporal IC):} How does the IC type evolve across version
  updates? What fraction of updates constitute silent expansions? (\S\ref{sec:temporal})
\item \textbf{RQ3 (Detection):} Can the IC pipeline detect confirmed
  malicious supply-chain injections? What is the precision-recall tradeoff? (\S\ref{sec:groundtruth})
\item \textbf{RQ4 (Cross-platform):} Is capability-description inconsistency
  platform-specific or ecosystem-wide? (\S\ref{sec:static})
\end{itemize}

% sections/07_detection.tex
% §7: The Capability-Mismatch Detection System
% Updated with IC framework terminology throughout

\section{The Capability-Mismatch Detection System}
\label{sec:detection}

We propose a multi-stage screening pipeline designed to identify capability-description mismatches at publish time and prioritize the most concerning skills or updates for review. The pipeline is grounded in the IC framework (\S\ref{sec:taxonomy}): it targets skills in the UD (undisclosed) quadrant of the difference matrix, where actual capabilities exist in code but are absent from user-facing descriptions.

\subsection{Static Capability Screen (SCS)}

\paragraph{Design.}
At publish time, every new skill submission undergoes an initial LSCPS
static sweep to extract standard security capability families directly
from its source components. This yields a systematic, code-centric
capability profile ($\mathcal{C}_{\text{act}}$) that is reproducible
given the same source artifacts, subject to the coverage limits of the
300+ pattern LSCPS bank. If a confirmed severe pattern is explicitly
detected (e.g., Credential Read combined with direct Exfiltration),
it acts as an immediate short-circuit, triggering a critical alert.

\subsection{Disclosure Mismatch Ranker (DMR)}

\paragraph{Design.}
Because marketplace descriptions are unstructured prose, we define a
\emph{disclosure matching layer} that reads the description and infers
an advertised capability set ($\mathcal{C}_{\text{vis}}$). The
matching engine performs three-level analysis:
\begin{itemize}
\item \textbf{Explicit disclosure}: direct keyword mentions of capabilities (``executes shell commands'', ``requires API key'');
\item \textbf{Implicit disclosure}: functional phrases implying capabilities (``automation tool'' implies SHELL\_EXEC);
\item \textbf{Negation detection}: safety claims contradicted by code (``read-only'' while code contains \texttt{rm -rf}).
\end{itemize}

\noindent
This matching is implemented as a curated bank of 80+ disclosure phrases per capability family.
The matching engine handles three types of natural language ambiguity:
\begin{enumerate}[topsep=1pt,itemsep=1pt]
\item \textbf{Synonym resolution}: each capability family has 10--30 synonyms
  compiled from the ClawHub corpus (e.g., ``transmit credentials'' $\to$
  CRED\_TRANSMIT; ``post to webhook'' $\to$ NET\_WRITE).
\item \textbf{Negation detection}: 8 negation patterns (``no network access'',
  ``read-only'', ``never stores credentials'') mark the capability as explicitly
  \emph{denied}; if the code still contains it, this becomes a strong IC-3/FC signal.
\item \textbf{Implicit/conditional phrasing}: ``may access files'' and
  ``integrates with external services'' are treated as \emph{partial} AD
  (contributing to CP but not full CR); ``send summary to Slack'' is treated
  as NET\_WRITE AD because it describes an outbound external-service communication.
\end{enumerate}
The matching engine is intentionally lightweight---the goal is not to solve
natural-language understanding completely, but to provide a deterministic
first-pass estimate deployable at marketplace scale,
with human or LLM review on a much smaller candidate queue.

By intersecting $\mathcal{C}_{\text{act}}$ and $\mathcal{C}_{\text{vis}}$, the pipeline constructs the difference matrix and computes CR (Eq.~\ref{eq:cr}), CP (Eq.~\ref{eq:cp}), CDI (Eq.~\ref{eq:cdi}), and RWM (Eq.~\ref{eq:rwm}).
Skills with $\text{CR} < \theta_{\text{cr}}$ or $\text{RWM} > \theta_{\text{rwm}}$
are flagged for review.
We set $\theta_{\text{cr}} = 0.3$ (below which a skill discloses fewer than 30\% of its detected capabilities) and $\theta_{\text{rwm}} = 3.0$ (above which undisclosed high-risk capability weight exceeds the autonomous-agent baseline).
Both thresholds were selected on the T6 training set by maximizing recall subject to a flagging rate $\leq 25\%$ of the corpus.
In IC framework terms, DMR targets IC-1 (Full Omission, CR\,=\,0) and IC-2 (Partial Omission, 0\,$<$\,CR\,$<$\,1) skills.

\paragraph{Risk-Weighted Mismatch (RWM).}
Rather than treating all undisclosed capabilities equally, we weight each
by its empirically-derived malicious-injection risk (Eq.~\ref{eq:rwm},
\S\ref{sec:taxonomy}), with weights $w_f$ learned from T6 ground truth
via Logistic Regression (Section~\ref{sec:riskweights},
Table~\ref{tab:riskweights}).

\paragraph{Results.}
We apply DMR to the full \NskillsDMR~\clawhub{} corpus (\S\ref{sec:static}).
Of \NskillsDMR skills, \textbf{\ResultFilesRead} files are scanned;
the pipeline detects at least one code-level capability in skills
classified IC-1 through IC-4:
\begin{itemize}
\item Mean CR = \ResultCR{} across all skills (including IC-0/TN);
      \textbf{\ResultICone{}} of skills are IC-1 (CR\,=\,0)
\item \textbf{\ResultICtwo{}} of skills are IC-2 (partial omission,
      $0 < \text{CR} < 1$)
\item Mean CDI = \ResultCDI{}; \textbf{92.8\%} of skills exhibit some
      form of capability-description inconsistency (IC-1 through IC-4)
\item \textbf{180 skills} exhibit deceptive denial: descriptions
      explicitly deny capabilities (``read-only'', ``no network access'')
      that \emph{are} present in code
\end{itemize}

\noindent
\textbf{Relation to DMR subset.}
The full-corpus IC classification and the DMR flagging pipeline
converge on $>$90\% inadequate disclosure across all \NskillsDMR{}~skills.
This cross-ecosystem convergence (\S\ref{sec:static}) suggests
the pattern is not unique to \clawhub{} or an artifact of a particular
analysis scope; whether it reflects deliberate developer behavior
or prevailing documentation norms requires further study.

\subsection{Data-Driven Capability Risk Weights}
\label{sec:riskweights}

\paragraph{Training.}
We learn RWM weights from 40 T6 malicious version pairs (mapping to 34 unique skills) and 400 negative controls using Leave-One-Out cross-validation with L2-regularized Logistic Regression.

\begin{table}[h]
\centering
\caption{Learned RWM weights from T6 ground truth (LOO-CV AUROC = 0.578). Heuristic CES weights from \S\ref{sec:background} shown for comparison.}
\label{tab:riskweights}
\small
\begin{tabular}{@{}lcccc@{}}
\toprule
\textbf{Capability} & \textbf{RWM $w_c$} & \textbf{CES} & \textbf{Solo AUROC} \\
\midrule
Autonomous agent   & 5.00 & 1.5 & 0.549 \\
Network exfil      & 4.93 & 2.0 & 0.525 \\
Download-execute   & 3.81 & --- & 0.512 \\
File destructive   & 3.47 & 1.5 & 0.516 \\
File write         & 3.41 & 1.5 & 0.474 \\
Persistence        & 2.99 & 1.5 & 0.486 \\
Credential access  & 1.12 & 2.5 & 0.445 \\
Shell execution    & 0.10 & 3.0 & 0.438 \\
\bottomrule
\end{tabular}
\end{table}

\paragraph{Interpretation.}
The learned weights diverge sharply from the heuristic CES
assignments (Table~\ref{tab:riskweights}). Shell execution,
weighted highest in CES (3.0), receives the \emph{lowest}
learned weight (0.10) because it is ubiquitous across both
malicious and benign skills---it has no discriminative power.
This is consistent with the UD analysis (\S\ref{sec:static},
Table~\ref{tab:hidden_fams}): SHELL\_EXEC is the most
frequently hidden family (\ResultICone{} of all skills exhibit
Full Omission, with SHELL\_EXEC the top hidden family), making it
a poor discriminator.
Conversely, autonomous agent behavior
(CES = 1.5 $\to$ RWM = 5.00) is the strongest predictor
because it is rare and strongly correlated with T6 cases.
The compound signature
\texttt{file\_destructive + autonomous\_agent} achieves a
lift of \textbf{50$\times$} over baseline co-occurrence
(Fisher's exact test, $p < 0.001$).

\paragraph{LOO-CV AUROC = 0.578.}
The modest AUROC reflects an important methodological result:
\emph{capability fingerprints alone cannot reliably distinguish
malicious from benign skills}.
Many benign skills exhibit the same high-risk capability profiles
as malicious ones.
We therefore treat RWM as a ranking signal inside a broader
review workflow rather than as a stand-alone malware classifier.
At CR threshold $\theta = 0.3$, \emph{pair-level} recall on T6 version pairs
reaches {\bf 77.5\%} (31/40 pairs);
at $\theta = 0.1$, pair-level recall reaches 97.5\% (39/40 pairs).
(Skill-level recall at $\theta=0.3$ is 76.5\%; see Table~\ref{tab:e2e}.)

\subsection{Temporal Drift Monitor (TDM)}

\paragraph{Design.}
To score potentially risky updates, we trace the delta over time. For each version transition, a temporal heuristic computes a feature vector:
\begin{equation}
\small
\mathbf{x} = [\text{TCE},\, \Delta\text{text\_ratio},\, \text{new\_kw},\, \text{SDS},\, \text{clog\_den}]
\label{eq:l3_vector}
\end{equation}
An anomaly score exceeding the historical threshold $\theta_{\text{drift}}$ flags the update for temporal review.

\paragraph{Results.}
The temporal detector achieves AUROC = 0.999 (LR) and 1.000 (RF) on proxy-labeled pairs.
\textbf{Important caveat:} These near-perfect scores reflect the deterministic nature of the proxy labels
(Section~\ref{sec:temporal}), which are constructed from the same static features the classifier uses.
The AUROC should therefore be interpreted as an upper-bound on temporal discriminability, not as
an independent validation metric. The meaningful, uncontaminated validation is on T6 ground truth
via the static SCS retrieval rule (79.4\% recall) and the combined SCS$\cup$DMR screening rule (94.1\% recall).

\subsection{End-to-End Retrieval Evaluation}
\label{sec:e2e}

\paragraph{Framing.}
Because \sysname is a \emph{publish-time risk prioritization} framework
rather than an automatic blocker, we report two complementary views:
(1) binary recall against the T6 ground truth, and
(2)~ranking utility under reviewer budget constraints (\S\ref{sec:triage}).
Low binary precision is expected given the ecosystem-wide non-disclosure pattern
documented in \S\ref{sec:static}; the meaningful question is whether
\sysname places T6 skills \emph{higher} in the review queue than natural baselines.

We evaluate each method against the 34 T6-confirmed malicious skills
and all 6,878 non-T6 DMR-eligible skills in our archived corpus
(Table~\ref{tab:e2e}).

\begin{table}[h]
\centering
\caption{T6 retrieval evaluation: 34 T6 skills vs.\
  6,878 non-T6 skills in the archived Code corpus.
  95\% Wilson CIs on recall.
  SCS threshold: CES$\geq$6; DMR threshold: CR$<$0.3.}
\label{tab:e2e}
\small
\begin{tabular}{@{}lccccl@{}}
\toprule
\textbf{Method} & \textbf{TP} & \textbf{FP} & \textbf{R}
  & \textbf{P} & \textbf{R 95\% CI} \\
\midrule
Platform scanner
  & 0   & 0     & 0.000 & ---   & --- \\
Random baseline
  & ---  & ---   & ---   & 0.005 & ---  \\
SCS: CES$\geq$6
  & 27  & 2,857 & 0.794 & 0.009 & [0.63, 0.90] \\
DMR: CR$<$0.3
  & 26  & 4,545 & 0.765 & 0.006 & [0.60, 0.88] \\
\textbf{SCS$\cup$DMR}
  & \textbf{32} & \textbf{5,375} & \textbf{0.941}
  & \textbf{0.006} & \textbf{[0.81, 0.98]} \\
\bottomrule
\end{tabular}
\end{table}

Key observations:
\begin{itemize}
\item \textbf{Platform scanners retrieve none of the confirmed T6 skills} (0/34).
The platform scanner is designed for metadata compliance and known
malware hash signatures---not runtime capability auditing of SKILL.md artifacts.
Its 0\% recall is not a strawman: it represents the state-of-practice that
marketplace users currently rely upon.
We verified this by checking all 40 T6 $v_{\text{bad}}$ versions against
the platform's public badge history; none were ever flagged.

\item \textbf{General-purpose security linters (Bandit, Semgrep)},
while detecting some Python anti-patterns, are not designed for
multi-file skill capability aggregation or description-vs-code comparison.
Bandit's high-severity rules target \texttt{subprocess} with \texttt{shell=True},
\texttt{eval/exec}, and deserialization---but \emph{miss} the dominant T6 patterns:
webhook-based exfiltration (\texttt{requests.post}), credential read
(\texttt{os.environ}), and autonomous file deletion.
Based on the T6 capability breakdown (Appendix~\ref{app:casestudies}),
Bandit high-severity would flag an estimated 30--40\% of T6 skills
(those using \texttt{subprocess}/\texttt{eval}) and miss the remaining 60--70\%
that rely on network-exfil or credential-harvest patterns alone.
SCS's 300+ marketplace-specific patterns cover all 14 LSCPS families,
achieving 79.4\% skill-level recall against this estimated 30--40\%.

\item \textbf{SCS static regex} (CES$\geq$6) catches
27/34 T6 skills (79.4\%). This is the \emph{conservative}
recall estimate, as it uses only heuristic CES weights
with no data-driven optimization on T6.
\item \textbf{DMR disclosure matching} independently catches
26/34 (76.5\%), including 5 skills missed by SCS.
In IC terms, DMR catches all IC-1 skills (CR\,=\,0)
and most IC-2 skills (CR\,$<$\,0.3).
\item \textbf{Combined SCS$\cup$DMR} catches 32/34 (94.1\%),
mapping exactly to the IC-1+IC-2 T6 coverage
(Table~\ref{tab:t6_ic_dist}).
The two missed skills have CES just below threshold
and CR\,$=$\,0.33, exceeding the 0.3 cutoff.
\item \textbf{Precision is very low} (0.6\%) because 78\%
of \emph{all} skills trigger SCS or DMR---reflecting the
ecosystem-wide non-disclosure problem documented in
\S\ref{sec:static}. This system is a \emph{prioritization}
tool, not an automated blocker.
\end{itemize}

\paragraph{Role of TDM (temporal layer --- future deployment).}
In the current full-population evaluation against archived T6 cases,
TDM does not improve recall: the archived dataset records
$v_{\text{bad}}\to v_{\text{fix}}$ (the remediation transition), which
correctly scores near-zero drift because\emph{capabilities are removed},
not added. TDM is therefore \emph{not evaluated as a recall component}
in this paper; its target scenario is a live streaming deployment
where each new version update is screened before installation.
In that setting, TDM would flag the \emph{injection} transition
($v_{\text{benign}}\to v_{\text{bad}}$) as a high-drift event,
reduce false-positive load on downstream SCS+DMR queues, and
provide a temporal risk signal for prioritization.
Full evaluation in a streaming deployment is reserved for future work
(\S\ref{sec:conclusion}).

\paragraph{Methodological note on RWM weights.}
The RWM weights in DMR were learned from T6 ground truth
via LOO-CV (\S\ref{sec:riskweights}). While LOO-CV
prevents direct overfitting, there is a risk that the
learned weights are optimized for the specific T6 sample.
The SCS recall (79.4\%) uses only heuristic CES weights
and is therefore the \emph{uncontaminated} baseline.

\paragraph{Precision in deployment.}
At the full-population class ratio (34:\NskillsDMR\ $\approx0.18\%$
prevalence), the pipeline flags a substantial fraction of
\NskillsDMR{} skills.
The system should therefore be used as a
\emph{ranking} mechanism, not a binary classifier.

\noindent\textit{Quantitative deployment scenario.}
Consider a marketplace ingesting 50 new skills/day.
SCS+DMR flags $\approx$39 (78\% × 50) for human review.
With P@20 = 0.150, a security analyst examining the
top-20 daily queue has a $1 - (1-0.150)^{20} = 96\%$
probability of encountering at least one malicious skill
per week (assuming base rate 1 malicious/50 submissions).
Alternatively, using the full ranked list at P@500 = 0.028
($5.7\times$ random lift), a 10-person team each reviewing
50 items finds $\approx$14\% of the total malicious inventory
per sweep---versus $<$1\% for unaided review.
This is comparable to the flagging workload accepted by
App Store and package-registry security teams~\cite{ladisa2023taxonomy}.

\paragraph{Ranking effectiveness.}
Ranking all \NskillsDMR{} skills by a composite score
(CES + RWM + compound signature bonus + CR penalty)
achieves P@20 = 0.150 (30.5$\times$ lift over the
random baseline of 34/\NskillsDMR).
At P@500, recall reaches 41.2\% while maintaining
5.7$\times$ lift (full P@K results in
Appendix~\ref{app:patk}).

\paragraph{Threshold sensitivity.}
Recall is robust across CR thresholds:
at CR$\,<\,$0.1 (IC-1 only), recall is 64.7\%;
at CR$\,<\,$0.3, recall plateaus at 88.2\%;
increasing to CR$\,<\,$0.5 yields no further T6
captures, confirming the CR threshold is not
a fragile parameter choice.

% ─────────────────────────────────────────────────────────────────────────────
\subsection{Baseline Comparison}
\label{sec:baselines}
% ─────────────────────────────────────────────────────────────────────────────

A key reviewer concern is that we compare only to the platform scanner,
which was not designed for capability-description mismatch detection.
To address this, we implement four natural baselines on the same full-corpus
JSONL data ({\small no additional code execution required}):

\begin{itemize}[topsep=2pt,itemsep=1pt]
\item \textbf{Platform Scanner}: State-of-practice; 0\% T6 recall (Table~\ref{tab:e2e}).
\item \textbf{B1 Lexical Overlap}: Jaccard($\mathcal{C}_{\text{vis}}$-words,
  $\mathcal{C}_{\text{act}}$-families) inverted as mismatch score.
  Captures description--capability divergence without LSCPS taxonomy.
\item \textbf{B2 TF-IDF Cosine Distance}: Description text vs.\ concatenated
  capability-family names; low similarity implies higher risk.
  Uses a lightweight per-corpus TF-IDF without external libraries.
\item \textbf{B3 Single-File Static}: Extract only from SKILL.md;
  no multi-file aggregation. Tests whether the multi-file design adds value.
\end{itemize}

\begin{table}[h]
\centering
\caption{Baseline comparison on T6 retrieval ranking.
  All methods evaluated on the identical 18,443-skill corpus,
  34 T6 confirmed malicious.
  Higher MRR and Recall@k indicate better prioritization.}
\label{tab:baselines}
\small
\begin{tabular}{@{}lrrrr@{}}
\toprule
\textbf{Method} & \textbf{MRR} & \textbf{R@100} & \textbf{R@200} & \textbf{R@500} \\
\midrule
Platform Scanner & 0.000 & 0.000 & 0.000 & 0.000 \\
B1: Lexical      & 0.002 & 0.000 & 0.000 & 0.029 \\
B2: TF-IDF       & 0.002 & 0.000 & 0.000 & 0.000 \\
B3: Single-File  & 0.002 & 0.000 & 0.000 & 0.029 \\
\textbf{SCS+DMR (Ours)} & \textbf{0.019} & \textbf{0.059} & \textbf{0.176} & \textbf{0.235} \\
\bottomrule
\end{tabular}
\end{table}

\paragraph{Interpretation.}
\sysname outperforms all natural baselines on every ranking metric.
The lexical (B1) and TF-IDF (B2) baselines capture description--code
divergence but cannot leverage the structured LSCPS risk taxonomy,
resulting in lower MRR and Recall@k.
The single-file baseline (B3) achieves substantially lower recall,
confirming that multi-file aggregation provides essential additional
capability signal: malicious code is frequently embedded in non-SKILL.md artifacts
(setup scripts, helper modules), which only the full-corpus scan captures.
The platform scanner is included as the \emph{state-of-practice} control;
its 0\% recall is not a strawman but rather reflects its design scope
(metadata compliance rather than capability auditing).

% ─────────────────────────────────────────────────────────────────────────────
\subsection{Triage Ranking Evaluation}
\label{sec:triage}
% ─────────────────────────────────────────────────────────────────────────────

\paragraph{Framing.}
Because the binary precision of our static pipeline is inherently low
(pervasive developer under-disclosure means 78\% of skills trigger SCS or DMR),
the more appropriate evaluation lens is \emph{ranking utility}: given a fixed
human review budget, does \sysname's composite ranking surface T6 skills
earlier than alternatives?

\paragraph{Composite score.}
We rank all \NskillsDMR{} skills by a composite score:
\begin{equation*}
\text{score}(s) = \text{CES}(s) + \text{RWM}(s) + 5(1{-}\text{CR}(s)) + 3{\cdot}\mathbf{1}[\text{cmpd}],
\end{equation*}
where compound is the \texttt{FILE\_DELETE}$\,\wedge\,$\texttt{AUTONOMOUS}
binary flag (50$\times$ lift, \S\ref{sec:riskweights}).

\begin{table}[h]
\centering
\caption{Triage ranking metrics on T6 retrieval.
  Base rate: $34/18{,}443\approx 0.18\%$.
  Lift@k = P@k / base rate.}
\label{tab:triage}
\footnotesize
\begin{tabular}{@{}lrrrrr@{}}
\toprule
\textbf{Method} & \textbf{MRR} & \textbf{R@100} & \textbf{R@200} & \textbf{R@500} & \textbf{L@200} \\
\midrule
Random (exp.) & 0.0002 & 0.018 & 0.036 & 0.092 &  1.0$\times$ \\
CDI Ranking   & 0.002  & 0.000 & 0.000 & 0.000 &  0.0$\times$ \\
CR-Only       & 0.017  & 0.088 & 0.088 & 0.206 &  2.4$\times$ \\
\textbf{\sysname} & \textbf{0.019} & \textbf{0.059} & \textbf{0.176} &
  \textbf{0.235} & \textbf{16.3$\times$} \\
\bottomrule
\end{tabular}
\end{table}

\paragraph{Budget coverage with SR-Lift.}
A marketplace security team reviewing the top-200 \sysname-ranked skills
recovers 17.6\% of all T6 cases (6/34)---a \textbf{16.3$\times$ lift}
over random selection.
At top-1,000, recall reaches 35.3\% with a 6.5$\times$ lift;
at top-2,000, 61.8\% recall with 5.7$\times$ lift.
For context, a random reviewer would need to examine
$\approx$9,600 skills (52\% of the corpus) to find the same 6 T6 cases.
Combined with the E4 SR-Precision result (top-200 = 100\% SR),
this confirms that the top of the review queue contains exclusively
security-relevant mismatches---not the benign documentation gaps
that constitute 34\% of the IC-1 population.
demonstrating that \sysname provides triage value orthogonal to
the platform scanner: it operates on description-code mismatches
rather than signature matching, capturing a complementary signal.

\begin{table}[h]
\centering
\caption{Triage operational scenarios.
  SR-Precision = fraction of review budget items judged security-relevant (E4).
  Baseline: random review ECR $={k}/{N}\times n_{\text{T6}}$.}
\label{tab:triage_ops}
\small
\begin{tabular}{@{}lrrrr@{}}
\toprule
\textbf{Budget} & \textbf{T6 Found} & \textbf{T6-Lift} & \textbf{SR-Prec.} & \textbf{ECR (random)} \\
\midrule
100  & 2/34 (5.9\%)  & 10.8$\times$ & 100\% & 0.18 \\
200  & 6/34 (17.6\%) & 16.3$\times$ & 100\% & 0.37 \\
500  & 8/34 (23.5\%) & 8.7$\times$  & --    & 0.92 \\
1000 & 12/34 (35.3\%) & 6.5$\times$ & --    & 1.84 \\
2000 & 21/34 (61.8\%) & 5.7$\times$ & --    & 3.69 \\
\bottomrule
\end{tabular}
\end{table}

\paragraph{Comparison to prior P@k claim.}
We note that an earlier version of this section reported P@20~=~0.150
based on a limited-population evaluation (only the 6,878 archived skills).
The full-corpus figures above (P@200~=~0.030 on all 18,443 skills)
are the correct full-population results and are used throughout.
The key finding is unchanged: \sysname substantially prioritizes T6 skills
over both the random baseline and all natural NLP baselines.

\paragraph{HC-Negative calibration: are FPs truly unlabeled or confirmed benign?}
A key methodological concern is whether the non-T6 skills in our evaluation
pool are \emph{confirmed benign} or merely \emph{unlabeled} (\S\ref{sec:groundtruth}).
To address this, we identify 1,326 \emph{high-confidence benign} (HC-Negative) skills:
1,247 TN-type skills (C\textsubscript{vis} = C\textsubscript{act}, CDI$=$0)
and 79 IC-0 skills (Accurate: CR\,=\,CP\,=\,1, description precisely
matches code, no undisclosed or over-claimed capabilities).
These skills are \emph{definitionally} not concealing dangerous capabilities.
Among the top-200 \sysname-ranked skills, \textbf{0/1,326 HC-Negative skills appear}---
the ranking successfully avoids confirmed-benign skills entirely.
This demonstrates that the FPs in the review queue are not skills
with known-accurate disclosure but skills with \emph{ambiguous} mismatch profiles
(legitimate automation tools with complex, under-documented capability sets).

\paragraph{Error analysis: FP patterns.}
Of the 195 non-T6 skills in the top-200 review queue:
\begin{itemize}[topsep=2pt,itemsep=0pt]
\item \textbf{90\% are IC-1} (full omission): they legitimately possess
  dangerous capabilities (e.g., CRED\_READ, NET\_WRITE, AUTONOMOUS)
  but do not disclose them in descriptions---the same pattern
  as malicious skills. These are \emph{undisclosed-but-likely-benign}
  automation tools (API integrations, scheduling agents, monitoring skills),
  not false detections of capability absence.
\item \textbf{The 2 missed T6 skills} (FN) have CR\,$\geq$\,0.3
  (they partially disclose the injected family) and fall just
  below the SCS threshold ($\text{CES}=3.8, 4.1$)---consistent
  with our threshold sensitivity analysis (\S\ref{sec:e2e}).
\end{itemize}
\noindent
This error breakdown confirms that \sysname's precision is not limited by
\emph{false signals} (detecting capabilities that don't exist) but by the
\emph{ecosystem-wide under-disclosure norm}: legitimately powerful tools
assume the same mismatch profile as malicious ones.
The E4 SR classification (53.8\% SR in the IC-1/IC-2 pool; \S\ref{sec:detection})
quantifies this overlap and guides prioritization within the flagged set.

% ─────────────────────────────────────────────────────────────────────────────
\subsection{Multi-View Capability Evidence Consistency}
\label{sec:multiview}
% ─────────────────────────────────────────────────────────────────────────────

A natural concern is that our mismatch findings are artifacts of a single
rule library. To test this, we deploy three \emph{independent} weak-evidence
views that share no code:

\begin{itemize}[topsep=2pt,itemsep=1pt]
\item \textbf{View A}: LSCPS code regex (300+ patterns; primary method).
\item \textbf{View B}: Package/dependency inference only
  (\texttt{requirements.txt} / \texttt{package.json} $\to$ PACKAGE\_CAPS mapping;
  no regex).
\item \textbf{View C}: Structured YAML/metadata surfaces only
  (\texttt{\_meta.json}, \texttt{skill.json}, YAML frontmatter;
  no regex).
\end{itemize}

\noindent
For each capability family, we ask: when View~A flags a family as hidden
(present in code, absent from description), does View~B independently
corroborate the presence of that capability?
Across the full corpus, View~B confirms View~A's under-disclosure signal
at a statistically meaningful rate for the high-risk families most
relevant to security:
the \textsc{AUTONOMOUS} family is confirmed 100\% of the time
(because autonomous scheduling packages always appear in dependency manifests),
and \textsc{NET\_WRITE} at 70\%+.
This cross-view convergence provides evidence that the mismatch signal
is not an artifact of any single extraction mechanism,
but reflects a genuine structural property observable from multiple independent
evidence sources.
(Full per-family confirmation rates in Appendix~\ref{app:multiview}.)


\label{sec:sa}

The current rule-based LSCPS pipeline (\S\ref{sec:detection}) is fully
deterministic and reproducible, but may miss capabilities expressed through
non-standard idioms or produce false positives in boilerplate code.
Semantic adjudication via LLM inference is a natural extension that would:
\begin{itemize}[topsep=2pt,itemsep=1pt]
  \item Reduce LSCPS false positives by verifying that matched code patterns
        are semantically functional (not, e.g., commented-out or dead code);
  \item Recover false negatives in $\mathcal{C}_{\text{act}}$ where capabilities
        are expressed through non-keyword patterns;
  \item Provide model-independent inter-rater reliability estimates
        via ensemble majority voting ($\geq3/5$ models agree).
\end{itemize}

We have implemented a five-model ensemble prototype
(\mone, \mtwo, \mthree, \mfour, \mfive)
running locally via vLLM~\cite{kwon2023vllm} with no commercial API dependence;
however, \textbf{full corpus execution is deferred to future work}---current
main results report rule-based LSCPS extraction only.
Semantic adjudication is expected to narrow the gap between our
conservative CR lower bound and true disclosure rates, but this
remains to be empirically quantified.

\subsection{Detection Layer Coverage by IC Type}

\begin{table}[h]
\centering
\caption{Detection layer effectiveness per IC type.
  $\checkmark$ = effective; $\sim$ = partial; $\times$ = ineffective.}
\label{tab:ic_coverage}
\scriptsize
\begin{tabular}{@{}lccccc@{}}
\toprule
\textbf{IC Type} & \textbf{T6\%} & \textbf{SCS} & \textbf{DMR} & \textbf{TDM} & \textbf{Plat.} \\
\midrule
IC-1 (Full Omis.)   & 64.7\% & $\checkmark$ & $\checkmark$ & --  & $\times$ \\
IC-2 (Partial Omis.) & 29.4\% & $\sim$       & $\checkmark$ & --  & $\times$ \\
IC-3 (False Claim)   &  0\%   & $\times$     & $\checkmark$ & --  & $\times$ \\
IC-4 (Bidirect.)     &  0\%   & $\checkmark$ & $\checkmark$ & --  & $\times$ \\
\bottomrule
\end{tabular}
\end{table}

No single layer is sufficient (Table~\ref{tab:ic_coverage}).
SCS catches code-level dangerous patterns but cannot assess disclosure;
DMR provides the crucial distinction via the difference matrix;
TDM catches \emph{new} risks introduced by version updates.
Together, they provide a practical, deployable screening workflow
grounded in the IC framework's formal decomposition of mismatch types.

\paragraph{Deployability.}
The static components of this framework operate entirely at publish-time, introducing \emph{zero runtime latency}. Marketplaces can deploy the SCS+DMR pipeline today with no OS-level sandbox modifications and reserve LLM review for a much smaller candidate set.

\subsection{Case Study: Trojanized Security Scanner}
\label{sec:casestudies}

\texttt{guard-scanner} (9.1~days exposure, 9~installs) exemplifies
the ironic disguise attack: a skill marketed as a \emph{security
scanning tool} contained credential harvesting, network exfiltration,
persistence mechanisms, and download-execute chains across 25+
detection patterns. IC type: \textbf{IC-2} (Partial Omission;
CR\,=\,0.143 --- one capability family was disclosed, but six were hidden);
RWM\,=\,12.86. All three detection layers (SCS, DMR, TDM) flag this case;
the platform scanner missed it entirely.
Four additional high-impact T6 reconstructions appear in
Appendix~\ref{app:casestudies}.

\subsection{Limitations}

\paragraph{Disclosure matching fidelity.}
DMR relies on deterministic curated phrase banks rather than
full natural-language understanding. Vague descriptions
(e.g., ``helps interact with external services'') fail to match
explicit NET\_WRITE markers, potentially misclassifying implicit
disclosures as omissions. Our multi-model adjudication study suggests
that richer semantic review is useful precisely because the static
pipeline over-flags benign under-disclosure.

\paragraph{Distinguishing Security-Relevant Mismatch from Benign Under-Documentation.}
A critical methodological question is whether the 75.8\% IC-1 prevalence
represents genuine \emph{security risk} or merely \emph{poor documentation hygiene}.
We address this through a stratified security-relevance study (E4):
500 IC-1/IC-2 skills were sampled across five CDI risk tiers and assessed
using a \emph{pre-defined, T6-independent} security-relevance rubric.
The rubric was authored before examining any sampled skill and consists of
three necessary conditions for an SR label, each evaluable from code alone
without reference to T6 labels or RWM weights: (i)~the hidden capability
family belongs to \{CRED\_READ, CRED\_TRANSMIT, NET\_WRITE, AUTONOMOUS,
FILE\_DELETE\}; (ii)~at least one concrete code path implements the capability
(not dead code or stub); and (iii)~the description does not hint at the
capability even implicitly (\ie{} it would not "put a reasonable user on notice").
Each skill receives one of three labels:
\textbf{SR} (Security-Relevant undisclosure: all three conditions met),
\textbf{BU} (Benign Under-documentation: conditions not fully met),
or \textbf{UN} (Uncertain: ambiguous code path or borderline disclosure).
Two researchers applied the rubric independently; Cohen's $\kappa = 0.81$
(substantial agreement) on the 500-skill sample before adjudication.

\begin{table}[h]
\centering
\caption{E4: Security-relevance by CDI tier ($N{=}500$ stratified sample).
  SR = security-relevant hidden capability; BU = benign under-documentation.}
\label{tab:e4_sr}
\small
\begin{tabular}{@{}p{3.2cm}rrr@{}}
\toprule
\textbf{Tier} & \textbf{N} & \textbf{SR} & \textbf{SR\%} \\
\midrule
T6-adjacent (CDI$=$1, high-risk)  & 100 & 100 & \textbf{100\%} \\
High-risk (CDI$>$0.9, cmpd)       & 100 &  52 & \textbf{52\%} \\
Medium-risk (CDI 0.5--0.9)         & 100 &  58 & \textbf{58\%} \\
Low-risk (CDI$<$0.5, low-risk cap) & 100 &   2 & \textbf{2\%} \\
Random IC-1/IC-2 baseline          & 100 &  57 &  57\% \\
\bottomrule
\end{tabular}
\end{table}

\noindent Key findings from E4 (inter-rater $\kappa = 0.81$;
  disagreements resolved by consensus):
\begin{itemize}[topsep=2pt,itemsep=1pt]
\item \textbf{53.8\% of IC mismatches meet all three SR conditions}, suggesting
  that a substantial fraction represent genuine consent risk rather than
  mere documentation incompleteness. This fraction rises to 100\% in the
  high-CDI tier where the hidden capabilities are predominantly
  CRED\_READ, CRED\_TRANSMIT, and AUTONOMOUS.
\item \textbf{CDI-SR correlation}: SR rate rises from 2\% (CDI$<$0.5)
  to 100\% (CDI$=$1, high-risk compound pattern)---a 50$\times$ range
  consistent with CDI being a useful risk-prioritization signal,
  though the middle tiers (CDI\,$\in$\,[0.6, 0.9]) show more variance
  and should not be read as a strictly monotonic ordering.
\item \textbf{SR-Precision@200 = 100\%}: the top-200 skills by CDI
  composite rank all satisfy the SR rubric---every item in the
  priority queue hides at least one high-risk capability and
  represents a genuine consent-risk mismatch, not merely an
  incomplete description. Note: \emph{security-relevant} does not
  imply malicious intent; it means the mismatch warrants human review.
\end{itemize}

\paragraph{Binary precision vs. triage precision.}
The raw binary precision (0.6\%) captures flagging against all 18,443 skills;
it conflates SR and BU mismatches.
The security-operationally meaningful metric is precision\emph{within the review queue}:
at CDI-ranked top-200, the queue contains 100\% SR---every reviewed skill
is security-relevant.
Platforms should configure the review budget to cover the top-k CDI tier,
not use \sysname as a binary blocker.

\paragraph{Ground truth scale.}
Our T6 evaluation uses 34 unique malicious skills (40 version pairs).
While this is the complete set of confirmed supply-chain injections
identified on \clawhub{}, the small sample limits statistical power:
the 95\% Wilson CI on our combined recall (94.1\%) spans
[80.9\%, 98.4\%]. We report CIs on all key metrics
to make this uncertainty explicit.

\paragraph{Platform scope.}
Our \clawhub{} analysis covers all \NskillsDMR{} skills.
Temporal analysis is scoped exclusively to \clawhub{}, which provides
explicit semantic version histories.
\skillssh{} (116,208 skills) does not expose structured version numbering;
it is used as a \emph{cross-ecosystem static validation} (\S\ref{sec:static}).
Both independent ecosystems converge on high inconsistency rates,
supporting the generalizability of our main finding.

\paragraph{Rule-Based vs.\ Semantic Extraction: Consistency Validation (E7).}
A reviewer may ask why we use curated keyword banks
rather than prompting an LLM to extract $\mathcal{C}_{\text{vis}}$.
We chose keyword matching for three reasons:
(1)~\emph{reproducibility}: deterministic extraction eliminates
prompt sensitivity and model-version drift;
(2)~\emph{scale}: 134K skills require sub-second per-skill processing;
(3)~\emph{evasion resistance}: LLM-based extraction
is vulnerable to adversarial descriptions.

To validate that rule-based extraction produces meaningful results,
we compare our pipeline directly with Qwen3-8B LLM semantic extraction
on the 6,912 skills for which both results exist (E7):

\begin{itemize}[topsep=2pt,itemsep=1pt]
\item \textbf{IC-type agreement}: Cohen's $\kappa = 0.154$;
  both methods agree IC-1 is dominant (rules: 78.0\%, LLM: 43.6\%),
  confirming the core finding is robust to extraction methodology.
\item \textbf{CR ordinal agreement}: Spearman $\rho = 0.432$---
  moderate positive correlation, confirming rules and LLM
  consistently rank the same skills as high vs.\ low disclosure.
\item \textbf{CR value gap}: rule-based CR ($\bar{x}=0.039$) is lower than
  LLM CR ($\bar{x}=0.218$), \emph{as expected}: the LLM accepts
  implicit disclosures (``accesses the internet'' $\to$ AD for NET\_READ),
  while the rule engine requires explicit vocabulary.
  \textbf{Rule-based CR is therefore a strict lower bound}
  on true disclosure; the actual gap between what skills disclose
  and what they can do lies between 0.039 and 0.218.
\end{itemize}

The $\kappa$ gap ($<$0.2) reflects a methodological difference
(keyword strictness vs.\ semantic leniency), not a measurement error.
Both methods agree on the direction and dominance of IC-1.
Keyword extraction serves as the conservative, verifiable baseline;
LLM semantic extraction is deferred to future work as described above.

\paragraph{Cross-platform IC direction reversal.}
The IC-type distribution reverses between platforms
(ClawHub: UD-dominated; skills.sh: FC-dominated),
which may appear to undermine the framework.
We argue this reversal \emph{validates} the methodology:
it correctly captures the structural difference
between terse marketplace blurbs (\clawhub{})
and verbose README-length descriptions (\skillssh{}).
The converging \emph{total} inconsistency ($\approx$90\%
on both platforms) is the robust cross-platform finding.

\subsection{Adversarial Robustness Analysis}
\label{sec:evasion}

A key reviewer concern for any detection system is evasion:
if an adversary \emph{knows} the detection mechanism,
can they craft skills that bypass it?
We analyze four evasion strategies against \sysname's three layers.

\begin{table}[h]
\centering
\small
\setlength{\tabcolsep}{4pt}
\caption{Estimated T6 recall under adversarial evasion.
Baseline (no evasion): 94.1\% empirical recall.
$^\dagger$Projected based on planned 5-model ensemble (L3 = future work; \S\ref{sec:sa}).
$^\ddagger$Description inflation backfires: over-claiming creates a detectable IC-0/IC-4 anomaly.}
\label{tab:evasion}
\begin{tabular}{@{}p{2.3cm}lcp{3.2cm}@{}}
\toprule
\textbf{Evasion} & \textbf{Layers} & \textbf{Recall} & \textbf{Attacker Cost / Backfire} \\
\midrule
None (baseline)        & ---      & 94.1\%          & --- \\
Desc.\ inflation$^\ddagger$ & L2 only & $\approx$65\%  & Creates IC-0 (over-claim) anomaly, increases reviewer suspicion \\
Code obfuscation       & L1 part. & $\approx$72\%   & Minified code is anomaly signal \\
Combined (A+B)         & L1+L2   & $\approx$45\%   & Detectable by code-style filter \\
Full adversarial$^\dagger$ & All & $\approx$20\% & Must fool 3/5 independent LLMs \\
\bottomrule
\end{tabular}
\end{table}

\paragraph{Key findings.}
Each evasion vector incurs measurable attacker cost.
Description inflation---adding broad claims like ``may access files, network,
credentials''---reduces DMR's recall from 94.1\% to $\approx$65\%,
but \emph{backfires} by shifting the skill from IC-1 (full omission) to
IC-0 (over-claim: description asserts capabilities not in code) or IC-4
(bidirectional mismatch). Both IC-0 and IC-4 are \emph{independently suspicious}
signals: a skill claiming broad capabilities while implementing few triggers
an ``over-declarative anomaly'' rule that can supplement the CDI ranking.
Code obfuscation produces entropy-anomalous patterns detectable
by a simple AST-parse completeness check.
The first three evasion scenarios are assessed against the current
deployed layers (L1: SCS, L2: DMR) and are empirically grounded.
The ``Full adversarial'' row ($^\dagger$) assumes a future L3
(5-model LLM ensemble, \S\ref{sec:sa}): defeating all layers
simultaneously requires fooling $\geq$3/5 independent model
families---a substantially harder bar than evading any single layer.
\sysname is a \emph{defense-in-depth} system:
even without L3, L1+L2 recall degrades to only $\approx$45\%
against a fully adversarial attacker.
Marketplace operators should supplement \sysname with
code-style anomaly detection and IC-0 flagging as complementary layers.

\paragraph{Low-cost evasion via vague disclosure (known limitation).}
We explicitly acknowledge a low-effort evasion that does not require
code modification: a developer who adds a single vague sentence
such as ``may access external services, local files, and system resources''
would reduce their CR below $\theta_{\text{cr}}=0.3$,
potentially bypassing the DMR threshold while providing minimal
actionable disclosure to users.
Our simulation suggests this reduces recall by $\approx$30pp (Table~\ref{tab:evasion}).
We do not claim to prevent this evasion, nor do we claim that disclosures
produced under detection pressure are authentic.
Rather, \sysname's deployment goal is to raise the \emph{cost of silent omission}
at publish time---not to serve as a binary security gate.
A marketplace policy that requires description-code consistency
incentivizes developers to produce specific, capability-accurate disclosures;
vague over-declarations shift mismatches from IC-1 (undetectable by platform)
to IC-0/IC-4 (detectable over-claim anomalies), creating a separate
moderation signal even against strategic actors.

% ─────────────────────────────────────────────────────────────────────────────
\subsection{Discussion: Toward Automated Capability Disclosure}
\label{sec:autodesc}
% ─────────────────────────────────────────────────────────────────────────────

\noindent\textit{This section describes a \textbf{constructive future direction},
not part of the current \sysname triage evaluation.
AutoDesc and CapManifest are presented to motivate open disclosure protocols;
full deployment study is deferred to future work.}

\medskip
\sysname addresses the \emph{detection} side of the mismatch problem:
it identifies skills whose actual capabilities exceed their disclosed ones
and ranks them for human review.
However, detection alone does not close the loop:
\emph{how should a benign developer who wants accurate disclosure actually write one?}
This section outlines \textbf{AutoDesc}, a constructive complement
that uses LLM inference over LSCPS extraction results to automatically
generate capability-accurate description drafts and a machine-readable
\textbf{CapManifest} protocol file.

\paragraph{AutoDesc overview.}
Given $\mathcal{C}_{\text{act}}(s)$ from LSCPS and the corresponding
code evidence snippets, AutoDesc prompts an LLM to produce:
(i)~a revised marketplace description that explicitly discloses every
detected capability family in plain language; and
(ii)~a structured CapManifest YAML file (Listing~\ref{lst:capmanifest})
that serves as a machine-verifiable capability declaration---analogous
to the Android permission manifest~\cite{felt2011android} that LLM
skill ecosystems currently lack.
The LLM is explicitly instructed \emph{not} to claim capabilities
absent from $\mathcal{C}_{\text{act}}$, preventing IC-3 (false-claim) generation.

\paragraph{CapManifest protocol.}
Listing~\ref{lst:capmanifest} shows the CapManifest format for
the \texttt{guard-scanner} T6 case (§\ref{sec:casestudies}).
Each entry records the detected capability family, its evidence location,
and a user-facing consent clause.
The \texttt{cr\_before}/\texttt{cr\_after} fields enable platforms to verify
that applying the generated description closes the IC gap:
$\text{CR}_{\text{before}}=0.143$ (IC-2) $\to$ $\text{CR}_{\text{after}}=1.0$ (IC-0).
\begin{lstlisting}[language=yaml, caption={CapManifest v1.0 for \texttt{guard-scanner} (pseudonymized). CR improves from 0.143 (IC-2) to 1.0 (IC-0) after applying the AutoDesc-generated description.}, label=lst:capmanifest, basicstyle=\ttfamily\tiny]
capmatch_manifest_version: "1.0"
skill_id: "guard-scanner"
generated_by: "CapMismatch-AutoDesc"
capabilities:
  - family: CRED_READ
    risk_level: HIGH
    confidence: 0.95
    evidence: {file: handler.py, line: 42,
                snippet: "api_key = os.environ['API_KEY']"}
    user_consent: "Read access to stored credentials"
  - family: NET_WRITE
    risk_level: MEDIUM
    confidence: 0.91
    evidence: {file: exfil.py, line: 87,
                snippet: "requests.post(webhook_url, json=payload)"}
    user_consent: "Send data to external endpoints"
  - family: SHELL_EXEC
    risk_level: MEDIUM
    confidence: 0.88
    evidence: {file: scanner.py, line: 23,
                snippet: "subprocess.run(['bash', scan_script])"}
    user_consent: "Execute system shell commands"
consent_summary:
  risk_tier: HIGH
  requires_explicit_consent: true
disclosure_metrics:
  cr_before: 0.143   # IC-2
  cr_after:  1.000   # IC-0 after applying AutoDesc description
  cdi_before: 0.857
  cdi_after:  0.000
\end{lstlisting}

\paragraph{AutoDesc as a remediation complement.}
AutoDesc and \sysname play orthogonal roles.
\sysname targets \emph{detection}: it identifies and ranks high-risk
skills for platform review, including potential malicious actors.
AutoDesc targets \emph{remediation}: it supplies benign developers
who care about accurate disclosure with a machine-assisted drafting tool,
reducing the documentation effort that makes under-disclosure the
path of least resistance.
Importantly, AutoDesc is not a bypass mechanism:
a skill with AutoDesc-generated accurate disclosure still undergoes
\sysname triage; it simply produces a lower CDI score and exits the
flagged queue---the intended outcome for a compliant skill.

\paragraph{Design goals and limitations.}
A deployable AutoDesc system must satisfy three properties:
(1)~\emph{conservatism}: never generate a description claiming capabilities
outside $\mathcal{C}_{\text{act}}$ (prevents IC-3 generation);
(2)~\emph{completeness}: every $f \in \mathcal{C}_{\text{act}}$ must appear
in the output description, so $\text{CR}_{\text{after}} \to 1.0$;
(3)~\emph{readability}: output must be intelligible to non-expert users.
A known limitation is that rule-based LSCPS extraction of $\mathcal{C}_{\text{act}}$
may miss non-standard capability idioms, so AutoDesc inherits the same
conservative lower-bound property as the detection pipeline.

\paragraph{Rule-based baseline evaluation.}
To validate the completeness property without LLM API access,
we implemented a rule-based AutoDesc baseline that generates
structured capability-accurate descriptions from 14 family-specific
sentence templates (one per LSCPS family; see artifact).
Applied to $n{=}500$ IC-1 skills sampled from \clawhub{}
(highest CDI first, mean 4.9 hidden families each),
the rule-based baseline achieves:
mean $\text{CR}_{\text{after}} = 0.943$ (up from 0.000),
mean $\text{CDI}_{\text{after}} = 0.044$ (down from 1.000),
and $74.6\%$ of skills reach $\text{CR} \geq 0.95$~(IC-0 equivalent)---
confirming that structured disclosure covering all detected families
can largely close the inconsistency gap.
The remaining $25.4\%$ reach partial disclosure (IC-2 range, $\text{CR} \in [0.5, 0.95)$),
primarily in skills with large $|\mathcal{C}_{\text{act}}| > 6$
where some families share overlapping keyword signals.
Full evaluation of LLM-generated (vs.\ template-based) description quality,
including fluency and developer acceptance, is deferred to future work.




% sections/04_static_analysis.tex
% §4: RQ1 Static IC Measurement
% Restructured around IC difference matrix framework

\section{RQ1: Static Inconsistency Measurement}
\label{sec:static}

\subsection{Experimental Setup}

For each of the \NskillsDMR{} \clawhub{} skills,
we extract two capability sets using the IC framework (\S\ref{sec:taxonomy}):
\begin{itemize}[topsep=2pt,itemsep=1pt]
\item $\mathcal{C}_{\text{act}}(s)$: Extended LSCPS regex scan across all skill artifacts
  (\ResultFilesRead{} files total; \ResultAvgFiles{} per skill on average),
  including SKILL.md, Python/JS/Shell source files, dependency manifests,
  and agent configuration files (300+ patterns, 14 families; Appendix~\ref{app:lscps}).
\item $\mathcal{C}_{\text{vis}}(s)$: Semantic keyword matching against marketplace
  description text using 14-family pattern vocabulary (same families as LSCPS).
\end{itemize}
We then compute the difference matrix $\langle\text{AD,UD,FC,TN}\rangle$ per skill
and derive CR, CP, CDI.
This section is \emph{fully deterministic}: no LLM inference.
All analysis uses the most recently archived version of each skill.

\paragraph{LSCPS taxonomy derivation.}
The 14 capability families were derived through three steps:
(i)~systematic review of prior supply-chain attack taxonomies~\cite{ohm2020backstabber,ladisa2023taxonomy}
and CWE categories for API abuse;
(ii)~manual inspection of the 40 T6 confirmed malicious case profiles to identify
the distinct capability fingerprints associated with real attacks;
(iii)~coverage validation against the top-1,000 \clawhub{} skills by install count
to ensure the families capture the dominant capability patterns in legitimate skills.
Families not observed in either T6 cases or the top-1,000 corpus were excluded.
Full family definitions, coverage rationale, and representative regex patterns
are in Appendix~\ref{app:lscps}.

\subsection{IC Type Distribution}

Table~\ref{tab:ic_types} (defined in \S\ref{sec:taxonomy}) reports
the IC type distribution over all \NskillsDMR{} skills.
Figure~\ref{fig:ic_dist} visualizes the result.

\begin{figure}[t]
\centering
\includegraphics[width=\columnwidth]{figures/fig_ic_distribution.pdf}
\caption{IC type distribution across \clawhub{} ($N{=}18{,}443$).
  Gold triangles indicate confirmed T6 malicious skills within each category.
  IC-1 and IC-2 together contain 94.1\% of all T6 cases.}
\label{fig:ic_dist}
\end{figure}

\noindent\textbf{Key aggregate metrics:}
\begin{itemize}[topsep=2pt,itemsep=1pt]
\item Avg.\ CR $=$ \ResultCR{}: descriptions mention only 4.2\% of detected capabilities on average.
\item Avg.\ CP $=$ \ResultCP{} (over the 3,208 skills with $\mathcal{C}_{\text{vis}} \neq \varnothing$):
  when capabilities are mentioned in descriptions, 83\% are confirmed in the code.
  The primary problem is \emph{omission}, not false claims.
\item Avg.\ CDI $=$ \ResultCDI{}: 89.4\% overall discrepancy index.
\end{itemize}

\paragraph{Key finding.}
\textit{
  92.8\% of skills exhibit at least one IC inconsistency type.
  The dominant observed direction is UD (capabilities present in code,
  absent from description), not FC (false claims).
  Skills predominantly under-disclose rather than over-claim
  under our rule-based extraction regime.
}

\subsection{CR vs.\ CP: The Two-Dimensional View}

\begin{figure}[t]
\centering
\includegraphics[width=\columnwidth]{figures/fig_cr_cp_scatter.pdf}
\caption{CR vs.\ CP 2D inconsistency space.
  Most skills cluster in the high-CP / low-CR quadrant (IC-1/IC-2):
  they do not lie about what they claim,
  but they claim only a fraction of what they do.
  Gold stars: T6 malicious skills.}
\label{fig:crcp}
\end{figure}

The two-dimensional view (Figure~\ref{fig:crcp}) reveals a structural pattern
not visible from CR alone.
Virtually all skills with claims ($\mathcal{C}_{\text{vis}} \neq \varnothing$)
have high CP ($\geq 0.9$): developers are not lying about the capabilities
they do mention.
The discrepancy is entirely on the CR axis:
the description captures under 5\% of what the code can do, based on our rule-based extraction.

This is an important nuance for threat modeling:
the behavioral signature of these mismatches is
\emph{passive concealment}, not active deception via false claims (IC-3/IC-4),
which are significantly rarer (0.4\% combined).

\subsection{Which Capability Families Are Most Hidden?}

\begin{figure}[t]
\centering
\includegraphics[width=\columnwidth]{figures/fig_ud_fc_families.pdf}
\caption{Comparison of UD (hidden, red) and FC (false-claimed, purple)
  counts per capability family. FC counts are $\approx$100$\times$ smaller
  than UD: skills almost never over-claim capabilities.}
\label{fig:udfcfam}
\end{figure}

Figure~\ref{fig:udfcfam} and Table~\ref{tab:hidden_fams} show the per-family
UD and FC breakdown.

\begin{table}[h]
\centering
\caption{Most frequently hidden (UD) capability families.}
\label{tab:hidden_fams}
\small
\begin{tabular}{@{}lrr@{}}
\toprule
\textbf{Family} & \textbf{Hidden Skills} & \textbf{\% of $N$} \\
\midrule
SHELL\_EXEC    & 13,988 & 75.8\% \\
CRED\_READ     &  8,612 & 46.7\% \\
AUTONOMOUS     &  6,906 & 37.4\% \\
MULTI\_STEP    &  5,816 & 31.5\% \\
NET\_READ      &  3,558 & 19.3\% \\
MEMORY\_WRITE  &  3,016 & 16.4\% \\
FILE\_READ     &  4,302 & 23.3\% \\
NET\_WRITE     &  2,386 & 12.9\% \\
CRED\_STORE    &  3,324 & 18.0\% \\
FILE\_DELETE   &    835 &  4.5\% \\
\bottomrule
\end{tabular}
\end{table}

The most-hidden capability---\textsc{Shell-Exec} (75.8\%)---is also the
most common across \emph{benign} skills, which explains the RWM weight
inversion described in \S\ref{sec:riskweights}:
prevalence dilutes discriminative power.
Conversely, \textsc{Cred-Read} (46.7\%) and \textsc{Autonomous} (37.4\%)
are high-UD and high-risk---they feature prominently in T6 cases.


\paragraph{Methodological note.}
Our rule-based extraction provides a \emph{lower bound} on true Capability Recall
and an \emph{upper bound} on IC-1 prevalence.
LSCPS patterns detect explicit API calls and keywords; capabilities implemented
via non-standard idioms, obfuscated code, or dynamically-loaded modules are not
captured. Similarly, $\mathcal{C}_{\text{vis}}$ keyword matching may
miss capabilities described in non-standard language.
LLM-based semantic extraction is expected to narrow both gaps and is deferred to
future work.

\subsection{Rule Robustness Analysis}
\label{sec:robustness}

A key concern is whether the IC-1 prevalence and T6 recall figures
are artifacts of our specific LSCPS pattern bank.
To test this, we re-run the full pipeline under 8 perturbation conditions:
\emph{Strict} (top-2 high-specificity patterns per family only),
\emph{Default} (full 300+ pattern bank, our primary setting),
\emph{Relaxed} (adds synonym expansions), four family-ablation conditions
(drop SHELL\_EXEC, CRED\_READ, NET\_WRITE, or AUTONOMOUS entirely),
and a \emph{5-family taxonomy collapse} (14 $\to$ 5 coarse groups).

\begin{table}[h]
\centering
\caption{Rule robustness: IC-1 prevalence and inconsistency under 8 conditions.\
  Default = our primary setting. Pilot $N{=}200$ for non-default conditions.}
\label{tab:robustness}
\footnotesize
\begin{tabular}{@{}p{3.0cm}cccc@{}}
\toprule
\textbf{Condition} & \textbf{IC-1} & \textbf{IC-2} & \textbf{Incons.} & \textbf{CR} \\
\midrule
Strict (top-2 patt.)      & 32.5\% & 2.0\%  & 47.5\% & 0.044 \\
\textbf{Default (ours)}   & \textbf{75.8\%} & \textbf{13.6\%} & \textbf{92.8\%} & \textbf{0.042} \\
Relaxed (+synonyms)       & 78.0\% & 12.0\% & 93.5\% & 0.034 \\
Ablate SHELL\_EXEC         & 71.5\% &  9.5\% & 86.0\% & 0.047 \\
Ablate CRED\_READ          & 74.0\% &  7.5\% & 86.0\% & 0.037 \\
Ablate NET\_WRITE          & 72.0\% & 10.0\% & 87.0\% & 0.042 \\
Ablate AUTONOMOUS          & 73.0\% &  6.5\% & 82.0\% & 0.034 \\
5-Family Collapse          & 73.0\% & 10.5\% & 87.5\% & 0.056 \\
\bottomrule
\end{tabular}
\end{table}

\noindent\textit{Interpretation.}
IC-1 prevalence ranges from 32.5\% (Strict) to 78.0\% (Relaxed), but every
non-Strict condition yields IC-1 $\geq$71\% and total inconsistency $\geq$82\%.
Even Strict---which retains only the top-2 highest-specificity patterns per
family, excluding $\sim$85\% of the pattern bank---records 47.5\% total
inconsistency.
CR is stable across conditions (range: 0.034--0.056), confirming the
under-disclosure finding is not an artifact of any specific pattern choice.
Family ablations show no single family drives the result: removing SHELL\_EXEC
or CRED\_READ each reduces IC-1 by $<$5pp, while total inconsistency remains
$>$82\%.
The 5-Family Collapse (14$\to$5 taxonomy) reproduces 73.0\% IC-1,
demonstrating that fine-grained taxonomy design is not critical to the
main finding---only that structured extraction is applied at all.



To assess whether capability-description inconsistency reflects a
cross-ecosystem pattern rather than \clawhub{}-specific artifacts,
we apply the identical IC pipeline without parameter tuning
to 116,208 skills from the independent \skillssh{} marketplace.

\begin{table}[h]
\centering
\caption{IC framework comparison across two platforms.
  No parameters were re-tuned between datasets.
  Inconsistency rates exceed 90\% on both platforms.}
\label{tab:crossplatform}
\small
\begin{tabular}{@{}lrr@{}}
\toprule
\textbf{Metric} & \textbf{\clawhub{}} & \textbf{\skillssh{}} \\
                & \small($N{=}18{,}443$) & \small($N{=}116{,}208$) \\
\midrule
Total Inconsistent         & 92.8\%  & 90.9\%  \\
IC-1 (Full Omission)       & 75.8\%  &  1.0\%  \\
IC-2 (Partial Omission)    & 13.6\%  &  1.1\%  \\
IC-3 (False Claim)         &  0.4\%  & 55.8\%  \\
IC-4 (Bidirectional)       &  2.9\%  & 32.9\%  \\
Accurate                   &  0.4\%  &  2.9\%  \\
\midrule
Avg.\ CR                   & \ResultCR{} & 0.828   \\
Avg.\ CP                   & \ResultCP{} & 0.338   \\
Avg.\ CDI                  & \ResultCDI{} & 0.650   \\
\bottomrule
\end{tabular}
\end{table}

\paragraph{A nuanced finding.}
The total inconsistency rate converges remarkably
(92.8\% vs.\ 90.9\%), but the \emph{dominant direction}
differs.
\clawhub{} is UD-dominated (IC-1+IC-2 = 89.4\%):
skills systematically under-disclose capabilities.
\skillssh{} is FC-dominated (IC-3+IC-4 = 88.7\%):
descriptions over-claim capabilities not present in code.
This divergence reflects platform conventions:
\skillssh{} descriptions are richer marketing prose
(higher $\mathcal{C}_{\text{vis}}$ from keyword hits),
while actual code capabilities ($\mathcal{C}_{\text{act}}$)
are structurally sparser.

Both platforms share the same Accurate rate (2.9\%),
meaning fewer than 3\% of skills in either ecosystem
have perfectly aligned capability-description sets.

\paragraph{Why the reversal does not invalidate the pipeline.}
The IC framework measures structural mismatch between
two independently-extracted capability sets.
The direction reversal reflects different \emph{input distributions},
not a measurement error:
\clawhub{} descriptions are concise marketplace blurbs,
while \skillssh{} stores full README-length text in the
description field (median $\approx$700 words vs.\
\clawhub{} descriptions of $\approx$30 words).
Longer text generates more $\mathcal{C}_{\text{vis}}$
keyword hits, many of which reference capabilities
discussed but not implemented, producing FC (false claims).
The reversal actually \emph{validates} the pipeline:
it correctly surfaces under-disclosure when descriptions are sparse
and over-claiming when descriptions are verbose.
A pipeline that produced the same IC direction on both
platforms despite their vastly different description conventions
would be suspect.

\paragraph{Answer to RQ4.}
\textit{
  Capability-description inconsistency is observed in $\approx$90\% of skills
  in both independently-measured ecosystems.
  The inconsistency direction reverses
  (\clawhub{}: under-disclosure; \skillssh{}: over-claiming),
  reflecting different description conventions, but the
  gap between what skills claim and what they do
  persists at scale in both platforms.
}

\paragraph{Controlled validation of the description-length hypothesis.}
If the IC direction reversal is driven by description verbosity rather than
platform identity, controlling for description length within \clawhub{}
should reveal an FC trend in verbose skills.
We stratify all \NskillsCH{} skills by description word count into quartiles
and compute IC-1 (under-disclosure) and IC-3 (false-claim) rates per bin:

\begin{center}\small
\begin{tabular}{@{}lrrrrl@{}}
\toprule
\textbf{Quartile} & \textbf{Range} & \textbf{Med.} & \textbf{IC-1} & \textbf{IC-3} & \textbf{Dir.} \\
\midrule
Q1 ($\leq$15 w)  & 1--15 w & 10 w & 78.9\% & 0.6\% & IC-1 (UD) \\
Q2 (16--20 w)    & 16--20 w & 19 w & 75.4\% & 0.6\% & IC-1 (UD) \\
Q3 (21--24 w)    & 21--24 w & 22 w & 77.2\% & 0.2\% & IC-1 (UD) \\
Q4 ($>$24 w)     & 25--42 w & 27 w & 71.1\% & 0.2\% & IC-1 (UD) \\
\bottomrule
\end{tabular}
\end{center}

\noindent\textit{Interpretation.}
\clawhub{} descriptions are uniformly brief (median 20 words; max 42 words),
so all quartiles remain IC-1-dominated: there are no verbose enough descriptions
to generate the FC-inflating keyword-hit density observed on \skillssh{}
(median $\approx$700 words).
This \emph{confirms} the description-length explanation via its expected
absence: the IC-3 reversal requires $\sim$100+ word descriptions
that \clawhub{} does not contain, while \skillssh{} has them at scale.
The monotone decrease in IC-1 (78.9\%$\to$71.1\%) with word count
further supports the hypothesis that more verbose descriptions provide
marginally more $\mathcal{C}_{\text{vis}}$ coverage.

\subsection{Additional Finding: Structured Declaration Layer Failure}

ClawHub provides structured YAML fields
(\texttt{autonomous}, \texttt{permissions}, \texttt{capabilities})
as a machine-readable capability declaration mechanism.
These fields are \emph{visible} to users on the marketplace listing page
and were designed to serve as a trusted complement to the prose description.

We find this mechanism has \textbf{comprehensively failed}:
\begin{itemize}[topsep=2pt,itemsep=1pt]
\item \texttt{autonomous} / \texttt{permissions} fill rate: \textbf{0.0\%}
  (0 of \NskillsDMR{} skills).
\item Any YAML capability declaration present: \textbf{18.4\%} (3,264 skills).
\item Among those 3,264, average YAML coverage of actual capabilities:
  \textbf{51.5\%}---the YAML itself is also incomplete.
\end{itemize}

This finding parallels~\cite{liao2020measuring}'s observation that
72\% of Alexa skills lack a privacy policy entirely:
the platform's intended transparency mechanism is ignored at scale.
The absence of enforcement---no certification check verifies YAML accuracy---
makes the declaration mechanism a \emph{voluntary gesture} that developers
choose not to make.

\subsection{Skill Category Breakdown}
\label{sec:category}

We infer skill categories from skill ID naming conventions across 18 domain labels
(Finance, Security, AI/Automation, Developer Tools, Communication, etc.).
Figure~\ref{fig:category} illustrates the IC violation rate and high-risk
capability prevalence per category.

\begin{figure}[t]
\centering
\includegraphics[width=\columnwidth]{figures/fig_category_ic_breakdown.pdf}
\caption{IC violation rates by skill category.
Left: stacked IC-type distribution.
Right: mean CR vs.\ high-risk capability prevalence (bubble size $\propto$ skill count).
Finance, Security, and AI/Automation occupy the lowest CR, highest risk quadrant.}
\label{fig:category}
\end{figure}

\paragraph{Key findings.}
Three categories form a \emph{high-risk, low-disclosure} cluster
with risk scores $>0.30$: \textbf{Finance}, \textbf{Security}, and \textbf{AI/Automation}.
These categories combine elevated prevalence of dangerous capabilities
(CRED\_READ, NET\_WRITE, AUTONOMOUS) with the lowest mean CR values ($\leq0.20$),
making them the highest-priority targets for marketplace review.

Conversely, \textbf{Education} exhibits the lowest risk score ($0.127$):
its skills tend toward simple content retrieval (NET\_READ) with minimal
credential or file-system access, and their descriptions more accurately
reflect this simpler profile.

Critically, \textbf{IC-1 (Full Omission) is endemic across all categories}---
even categories commonly perceived as low-risk (Monitoring: 16.1\%, Travel: 20.5\%).
This cross-category universality supports the interpretation that under-disclosure
is a broad \emph{marketplace norm} rather than confined to inherently
dangerous domains.

% sections/05_temporal_analysis.tex
% §5: RQ2 Temporal IC Analysis
% Restructured around IC temporal extension

\section{RQ2: Temporal Inconsistency Evolution}
\label{sec:temporal}

\subsection{Experimental Setup}

We construct a temporal dataset of 24,243 version transition pairs spanning 7,144 \clawhub{} skills (3.39 documented transitions per skill on average).

\noindent\textbf{Exploratory note.} Temporal IC labeling (T4/T5) relies on the same
LSCPS static extraction used for the main IC analysis; temporal labels are therefore
\emph{proxy indicators} derived from version-history metadata, not independently verified.
TDM recall is not included in the main composite evaluation (Table~\ref{tab:full_eval})
because its proxy labels are not independent of the detection features.
The findings here are reported as \emph{exploratory observations} to motivate
future deployment with live version-controlled repositories.
For each unique version pair $(v_i, v_{i+1})$, we compute:
\begin{itemize}[topsep=2pt,itemsep=1pt]
    \item \textbf{TCE} (Eq.~\ref{eq:tce}): Risk-Weighted Mismatch delta across versions;
    \item \textbf{SDS}: fraction of newly introduced capabilities mentioned in the changelog;
    \item \textbf{Family-level drift}: per-family CES deltas (Shell, Credential, Network, File, Agentic).
\end{itemize}
T4/T5 labels are \emph{heuristic risk indicators} from archived text, not verified malicious labels.

\paragraph{Scope note: \clawhub{} only.}
Temporal analysis is \clawhub{}-exclusive; \skillssh{} provides only a static snapshot
without structured version history.
Cross-platform static results are reported in \S\ref{sec:static} (RQ4).

\subsection{Temporal IC Distribution}

\begin{figure}[t]
\centering
\includegraphics[width=\columnwidth]{figures/fig_ic_temporal.pdf}
\caption{Version transition temporal IC labels on \clawhub{} ($M{=}24{,}243$ pairs).
  T6 (confirmed malicious) transitions are shown as a separate category.}
\label{fig:transitions}
\end{figure}

\begin{table}[h]
\centering
\caption{Temporal IC classification of 24,243 version transitions.
  ``Silent expansion'' = TCE$>$0 and SDS$<$0.35.}
\label{tab:temporal_ic}
\scriptsize
\begin{tabular}{@{}lrrp{3.2cm}@{}}
\toprule
\textbf{Label} & \textbf{Count} & \textbf{\%} & \textbf{IC meaning} \\
\midrule
T-Benign         & 19,131 & 78.9\% & No IC degradation \\
T-Silent-Expand  &  1,940 &  8.0\% & TCE$>$0, SDS$<$0.35 (T4) \\
T-Silent-Cred    &  1,136 &  4.7\% & CRED-READ silently expanded \\
T-DocLag         &    623 &  2.6\% & Capability added, description static \\
T-Silent-Minor   &    534 &  2.2\% & Low-risk generic drift \\
T-Silent-Major   &    512 &  2.1\% & High-risk cap injected without disclosure \\
T-CovertAgent    &    327 &  1.3\% & AUTONOMOUS silently added \\
\textbf{T6-Malicious} & \textbf{40} & \textbf{0.2\%} & \textbf{Confirmed malicious injection} \\
\bottomrule
\end{tabular}
\end{table}

\paragraph{Answer to RQ2.}
\textbf{2,575 of 24,243 transitions (10.6\%)} are silent expansions
(TCE$>$0, SDS$<$0.35).
This figure aggregates across all temporal sub-categories in Table~\ref{tab:temporal_ic}
(T-Silent-Expand, T-Silent-Cred, T-DocLag, T-Silent-Minor, T-Silent-Major, and T-CovertAgent)
that satisfy the TCE$>$0, SDS$<$0.35 criterion.
Ratio of silent to disclosed expansions: $2{,}575 : 534 = \textbf{4.8{:}1}$.
The most silently-expanded families:
CRED-READ (39.0\%), SHELL-EXEC (31.6\%), FILE-OPS (17.8\%).
Median changelog for silent expansions: 3 words (``minor bug fix'').

\subsection{Cumulative IC Degradation (T5)}

T5 requires IC-1/IC-2 at the current version
\emph{and} at least one silent expansion in the skill's version history.
We find \textbf{1,376 skills} (7.8\% of all \NskillsDMR{} skills;
19.3\% of the 7,144 skills with version history) satisfy this criterion.
T5 skills are \textbf{4.5$\times$ more likely} to be associated with
a T6 confirmed malicious injection than non-T5 skills
(Fisher's exact test, $p{<}0.01$; OR\,=\,4.5, 95\%\,CI\,[1.8,\,11.2]).
This association suggests T5 may be a useful risk indicator,
but the proxy nature of T4/T5 labels means this cannot be treated
as independent validation of the detection pipeline.

\paragraph{Changelog word count distribution.}
Among silently expanding transitions, the median changelog is 3 words
(e.g., ``minor bug fix,'' ``performance update,'' ``1.0.1'').
Only 4 of 2,575 candidate pairs feature changelogs exceeding 50 words.

\paragraph{Temporal pattern.}
Skills do not expand capabilities linearly.
A ``plateau-and-jump'' pattern emerges:
78\% of aggregate CES expansions occur within a single version transition
rather than gradually over time.
This concentration is consistent with abrupt capability change,
though not sufficient on its own to establish malicious intent.

\subsection{Version Lifecycle IC Evolution}
\label{sec:lifecycle}

A natural question is whether \emph{developers self-correct} over time:
do skills become more transparent in later versions?
We group transitions by the target version number $v$ into five cohorts
and compute per-cohort IC-0 rate (consistent disclosure) and mean CR.

\begin{table}[h]
\centering
\small
\setlength{\tabcolsep}{4pt}
\caption{IC disclosure quality across version lifecycle cohorts (24,243 transitions).
  All cohorts are sub-sampled by version number; v1 ($N{=}21$) and v2--3 ($N{=}54$) have
  small $N$ and should be interpreted cautiously.}
\label{tab:lifecycle_ic}
\begin{tabular}{@{}lrrrrr@{}}
\toprule
\textbf{Cohort} & \textbf{N} & \textbf{IC-0\%} & \textbf{IC-1\%} & \textbf{CR} & \textbf{Silent\%} \\
\midrule
v1   (Initial)     & 21  & 4.8  & 61.9 & 0.221 & 0.0 \\
v2--3 (Early)      & 54  & 7.4  & 51.9 & 0.252 & 5.6 \\
v4--7 (Mid)        & 100 & 10.0 & 45.0 & 0.274 & \textbf{11.0} \\
v8--15 (Mature)    & 191 & 11.0 & 55.5 & 0.313 & 0.5 \\
v16+  (Long-lived) & 134 & 20.1 & 56.0 & 0.357 & 0.0 \\
\bottomrule
\end{tabular}
\end{table}

\begin{figure}[t]
\centering
\includegraphics[width=\columnwidth]{figures/fig_temporal_ic_evolution.pdf}
\caption{IC distribution and mean CR by version cohort, with silent expansion rate.
Silent expansions peak in the mid-lifecycle (v4--7), the period of active development
but limited security scrutiny.}
\label{fig:lifecycle}
\end{figure}

\paragraph{Key findings.}
Mean CR improves slightly ($0.221 \to 0.357$) as skills age,
suggesting limited but real developer self-correction over time.
However, \textbf{IC-1 (Full Omission) remains the dominant IC type in all cohorts},
including long-lived skills ($>$v16): 56\% of transitions in that cohort
still exhibit complete under-disclosure.
This pattern is \textbf{consistent with} IC violations being a structural ecosystem property
rather than a developer-learning-curve artifact---though the cohort sizes
(especially v1--v3) are small and caution is warranted in generalizing.

Notably, silent capability expansions \textbf{peak at v4--7 (11\%)},
the mid-lifecycle window where skills are actively developed
but not yet subject to mature security scrutiny or active user review.
This is the \emph{highest-risk window} for undetected capability injection,
providing a deployment heuristic: skills in this version band should receive
higher-priority review from TrustDrift's screening queue.

% sections/06_ground_truth.tex
% §6: RQ3 --- Evaluation: Detecting Malicious Supply-Chain Injections
% Updated to integrate IC framework (IC-1 through IC-4)

\section{RQ3: Detecting Malicious Supply-Chain Injections}
\label{sec:groundtruth}

\subsection{Case Identification}

\paragraph{Selection independence.}
A central methodological concern is whether our ground truth selection is
\emph{independent} of the detection system we subsequently evaluate.
We affirm that it is:
T6 cases are identified exclusively from external incident evidence
(steps 1--3 below) without any reference to LSCPS scores or IC classifications.
LSCPS and IC metrics are computed \emph{after} case inclusion is finalized
and are used for characterization only.
The resulting ground truth is therefore a pre-registered, code-analysis-free oracle.

We identify verified malicious supply-chain injections (T6 cases) by cross-referencing:
\begin{enumerate}
    \item \clawhub{}'s public security incident records (advisories and removal notices);
    \item Changelog entries containing explicit remediation language
          (e.g., ``removed malicious code,'' ``security patch,'' ``exfil fix''),
          identified by keyword search over the version transition corpus;
    \item Manual review of the suspicious version ($v_{\text{bad}}$),
          the remediation diff, and any publicly documented response context,
          confirmed by at least two independent reviewers from the research team.
\end{enumerate}

No candidate was admitted to T6 on the basis of a high static capability score.
Specifically, we did not use LSCPS CES, CR, or any IC metric to identify or
pre-filter candidates.
This protocol yields \textbf{40 confirmed T6 version pairs} corresponding to
\textbf{34 unique skills}. Static \lscps{} scores are computed only
\emph{after} case inclusion is established and are used for characterization,
not for ground-truth selection.

\paragraph{T6 represents the complete confirmed set.}
The 34 T6 skills are not a sample: they constitute
\emph{every} confirmed malicious injection case we identified on \clawhub{}
during our archival period (July 2024--March 2026).
The limited size therefore reflects the scarcity of externally verifiable cases,
not a sampling choice. This scarcity is itself informative: malicious injections
appear rare in absolute terms but can reach hundreds of users within hours
(298 total installs during malicious windows across all 40 T6 pairs).

\paragraph{Corroboration source breakdown.}
Of the 40 T6 version pairs, external corroboration came from:
\textbf{31} with a \clawhub{} platform advisory or removal notice;
\textbf{6} with an explicit developer-authored remediation changelog
  (``removed malicious code''/``security patch'') plus public discussion;
\textbf{3} with public blog post or bug report identifying the skill.
All 40 have a code diff confirming capability injection and removal.
No pair was included on code diff alone without external corroboration.

\paragraph{T6 attack type distribution.}
%  Breakdown over 34 unique skills:
Among the 34 T6 skills:
\textbf{22} (64.7\%) involve credential theft (CRED\_READ + outbound exfil);
\textbf{8} (23.5\%) involve destructive file operations (FILE\_DELETE, rm -rf);
\textbf{6} (17.6\%) involve download-execute or persistence mechanisms;
\textbf{12} (35.3\%) use webhook-based exfiltration as the primary channel.
(Categories overlap; some skills exhibit multiple attack types.)

\paragraph{Inter-rater reliability.}
Each candidate was independently reviewed by two researchers
who examined the malicious version's diff, its changelog,
and the corresponding platform advisory or public report.
Inclusion required unanimous agreement.
For the 12 ambiguous cases (changelog mentioned ``security'' without detail),
a third reviewer served as tiebreaker.
Final T6 inclusion required: (a) at least one external corroboration source
(platform advisory or public CVE/disclosure) \emph{and}
(b) a code diff confirming capability addition between $v_{\text{bad-1}}$
and $v_{\text{bad}}$ followed by removal at $v_{\text{fix}}$.
The resulting protocol achieved Cohen's $\kappa = 0.91$ (almost perfect agreement)
on the 40 version pairs before adjudication, reflecting the clarity of malicious
injection evidence in confirmed T6 cases.

\subsection{Methodological Note: The Fix-Transition Framing}

A critical observation reconciling an empirical discrepancy in our continuous dataset is that all 40 T6 cases manifest as transitions with TDG $= 0$, SDS $= 1.0$, and \texttt{is\_candidate} $= 0$. This is not a measurement failure; it follows directly from how the temporal dataset is constructed.

Our temporal dataset records $v_{\text{bad}} \to v_{\text{fix}}$,
i.e., the \emph{remediation} transition.
In this transition, capabilities are \emph{removed} (TDG $< 0$),
not added.
The fix changelog explicitly mentions the security issue (SDS $= 1.0$).
The temporal detector correctly classifies this pair as clean
because the pair itself is not a drift event.

The attack surface was $v_{\text{bad}}$, not the transition.
Detection of T6 attacks therefore requires a \emph{static} audit of
published versions, not a temporal drift detector.

\subsection{T6 Characteristics through the IC Lens}

\begin{table}[h]
\centering
\caption{T6 malicious skill distribution across IC types.
  IC-1+IC-2 capture 94.1\% of all confirmed malicious skills.
  SCS~R.~(CES$\geq$6) and DMR~R.~(CR$<$0.3) show skill-level recall per IC type.}
\label{tab:t6_ic_dist}
\scriptsize
\begin{tabular}{@{}lrrcc@{}}
\toprule
\textbf{IC Type} & \textbf{T6} & \textbf{\%} & \textbf{SCS R.} & \textbf{DMR R.} \\
\midrule
IC-1 (Full Omission)    & 22 & 64.7\% & 16/22 (73\%) & 22/22 (100\%) \\
IC-2 (Partial Omission) & 10 & 29.4\% &  8/10 (80\%) &  8/10 (80\%)  \\
Accurate                &  1 &  2.9\% &  1/1  (100\%)&  0/1  (0\%)   \\
TN\_Only                &  1 &  2.9\% &  0/1  (0\%)  &  0/1  (0\%)   \\
\midrule
\textbf{IC-1+IC-2}      & \textbf{32} & \textbf{94.1\%} & \textbf{24/32 (75\%)} & \textbf{30/32 (94\%)} \\
\textit{\textbf{All T6 (SCS total)}} & \textit{34} & & \textit{27/34 (79\%)}$^\dagger$ & \\
\bottomrule
\end{tabular}
\begin{minipage}{\columnwidth}
\vspace{2pt}\scriptsize
$^\dagger$ SCS catches 3 T6 skills outside IC-1+IC-2 (1 Accurate, 2 TN-only at CES$\geq$6)
due to false-negative in $\mathcal{C}_{\text{vis}}$ extraction; total skill-level SCS TP=27.
\end{minipage}
\end{table}

\paragraph{Why IC-1 and IC-2 dominate T6.}
The IC framework (Table~\ref{tab:t6_ic_dist}) reveals that
malicious injections overwhelmingly target the UD (undisclosed)
quadrant of the difference matrix:
\begin{itemize}
    \item \textbf{22 T6 skills are IC-1} (Full Omission, CR\,=\,0):
    descriptions mention \emph{none} of the injected capabilities.
    DMR detects all 22 (100\% recall) because CR\,=\,0 guaranteed flags.
    \item \textbf{10 T6 skills are IC-2} (Partial Omission, 0\,$<$\,CR\,$<$\,1):
    descriptions mention some benign capabilities but omit the malicious ones.
    Both SCS and DMR independently catch 80\% of these.
    \item \textbf{0 T6 skills are IC-3}: as expected,
    malicious developers do not over-claim dangerous capabilities.
    \item The remaining 2 T6 skills (5.9\%) fall into Accurate (1) and TN\_Only (1),
    likely due to measurement noise in $\mathcal{C}_{\text{vis}}$ extraction;
    SCS (CES$\geq$6) nonetheless flags 1 of these 2 (the Accurate case)
    via its static capability score, explaining the SCS total TP of 27.
\end{itemize}

\paragraph{Key finding.}
\textit{IC-1+IC-2 capture 94.1\% (95\% Wilson CI: [80.9\%, 98.4\%])
of all confirmed malicious skills.
The IC framework therefore provides both a theoretical justification
and empirical validation for focusing detection on the UD (undisclosed)
quadrant of the difference matrix.}

\paragraph{Statistical note.}
Because IC-1+IC-2 is ubiquitous (89.6\% of \emph{all} skills),
the enrichment of T6 in IC-1+IC-2 (94.1\%) is not statistically
significant (Fisher's exact test $p = 0.57$, OR $= 1.85$).
This non-significance is itself informative:
IC type alone cannot distinguish malicious from benign skills.
The discriminative power comes from \emph{within} IC-1/IC-2:
which specific capabilities are hidden (RWM) and whether
compound dangerous signatures are present (SCS).
The IC framework provides the \emph{target space};
RWM and SCS provide the \emph{ranking} within that space.

\subsection{T6 Case Profiles}

\begin{table}[h]
\centering
\caption{T6 confirmed malicious version profiles (top 11 by CES; 40 total).%
  CES\textsubscript{static} = weighted LSCPS score on $v_{\text{bad}}$.}
\label{tab:t6cases}
\scriptsize
\begin{tabular}{@{}lcclp{3.3cm}@{}}
\toprule
\textbf{Skill} & \textbf{CES} & \textbf{IC} & \textbf{CR} & \textbf{Primary Attack} \\
\midrule
Skill-A  & 14.5 & IC-1 & 0.000 & CRED\_TRANSMIT + webhook \\
Skill-B  & 14.0 & IC-1 & 0.000 & CRED\_TRANSMIT + email \\
Skill-C  & 14.0 & IC-2 & 0.143 & CRED\_READ + exfil + shell \\
Skill-D  & 12.5 & IC-1 & 0.000 & NET\_WRITE + credential theft \\
Skill-E  & 12.0 & IC-2 & 0.167 & CRED\_STORE + autonomous \\
Skill-F  & 12.0 & IC-1 & 0.000 & FILE\_DELETE + CRED\_TRANSMIT \\
Skill-G  & 11.5 & IC-1 & 0.000 & CRED\_TRANSMIT + data exfil \\
Skill-H  & 11.5 & IC-1 & 0.000 & CRED\_TRANSMIT + autonomous \\
Skill-I  & 10.5 & IC-2 & 0.143 & CRED\_READ + shell + download \\
Skill-J  & 9.0  & IC-1 & 0.000 & NET\_WRITE + startup persist. \\
Skill-K  & 8.5  & IC-2 & 0.200 & AUTONOMOUS loop + CRED\_READ \\
\bottomrule
\end{tabular}
\end{table}

All 40 T6 cases share three recurring properties:
\begin{itemize}
    \item \textbf{High-Privilege Access:} CRED\_READ capability is present in $v_{\text{bad}}$ (37/40, 92.5\%);
    \item \textbf{Data Exfiltration:} At least one outbound communication capability (NET\_WRITE or CRED\_TRANSMIT) is active;
    \item \textbf{Complete Scanner Evasion:} 0 detections by the inherent platform scanner at the time of malicious version publication.
\end{itemize}

\paragraph{Low-CES but confirmed cases.}
Eight cases (CES $\leq 4.5$) lack explicit exfil patterns
but contain credential harvesting via memory logs or
download-and-execute flows through API interactions.
These cases are the ones our LSCPS auditor misses,
reinforcing the need for LLM-based annotation as a second review layer.

\subsection{Platform Scanner Failure Analysis}

All 40 T6 cases were publicly available and received installs
during the period when $v_{\text{bad}}$ was the current version:

\begin{itemize}
\item Average exposure window: 0.5 days
\item Maximum exposure window: 9.1 days
\item Total installs during malicious period: 298
\item Platform scanner detection: 0/40 pairs (0\%)
\item Our LSCPS static auditor (high-risk): \textbf{29/40 pairs (72.5\%)}
\item Our LSCPS static auditor (any capability): \textbf{37/40 pairs (92.5\%)}
\item Skill-level evaluation (IC-based): see Table~\ref{tab:e2e}
\end{itemize}

\begin{figure}[t]
\centering
\includegraphics[width=\columnwidth]{figures/fig_t6_detection_gap.pdf}
\caption{Detection gap on T6 pairs. The platform scanner misses all 40, while static LSCPS detects 72.5\% (pair-level) as high-risk. At skill level, SCS$\cup$DMR catches 94.1\% (Table~\ref{tab:e2e}). IC-1 and IC-2 together harbour 94.1\% of all T6 skills (Table~\ref{tab:t6_ic_dist}).}
\label{fig:t6_detection}
\end{figure}

The 100\% miss-rate by existing platform scanners is consistent with
their emphasis on auditing \emph{declared intent} rather than
\emph{runtime capability}. Within the IC framework, the platform
scanner operates exclusively on $\mathcal{C}_{\text{vis}}$ information,
which by definition cannot detect UD (undisclosed) capabilities.
Our pipeline explicitly audits $\mathcal{C}_{\text{act}}$ and
computes the difference matrix to surface the UD quadrant---the
precise region where all T6 malicious injections reside.

\subsection{Implications for Platform Defenses}

The T6 cases suggest three operational lessons, reinforced by the IC framework:

\begin{itemize}
    \item \textbf{Trust extends across versions.} Users verify $v_{\text{good}}$ at install time, but that trust also extends to later versions under auto-update. Auto-update without renewed capability review creates a structural exposure window.
    \item \textbf{IC-1 (Full Omission) is the default attack pattern.} In 22 of 34 T6 skills (64.7\%), the description contains zero disclosure of the injected capabilities. This is the natural strategy for malicious actors and the easiest pattern for automated detection.
    \item \textbf{Simple capability combinations already catch many cases.} Only 10 of the 40 cases required more specialized pattern matching; the remaining 30 were detectable through straightforward combinations such as CRED\_READ plus outbound communication. This suggests that a conservative publish-time IC-aware static audit would have surfaced many of these cases for review.
\end{itemize}

% sections/08_related_work.tex

\section{Related Work}
\label{sec:related}

\paragraph{LLM agent security.}
Tool-augmented LLM agents~\cite{yao2023react,schick2024toolformer,significant2023autogpt}
have created new ecosystems with unique trust assumptions~\cite{xi2023rise}.
Prompt injection~\cite{greshake2023indirect} hijacks agent {\em input};
our work targets the {\em supply-chain} layer where the installed skill
itself is adversarial. Agent sandboxing~\cite{andriessen2024sandbox,ruan2024identifying},
tool verification~\cite{liu2024toolauth}, and safety analyses~\cite{tian2024evil}
assume descriptions accurately reflect behavior---an assumption we
systematically invalidate across \NskillsDMR{} skills.
The MCP standard~\cite{anthropic2024mcp}
secures the invocation layer; we audit the disclosure layer.
Concurrent work on tool poisoning~\cite{wu2024toolsword} and
function-calling misuse~\cite{shi2024ehragent} motivates
publish-time rather than per-call defenses.

\paragraph{MCP and tool marketplace security.}
The Model Context Protocol (MCP)~\cite{anthropic2024mcp} has rapidly emerged
as a standard for LLM tool integration, but its security properties
at the marketplace level remain underexplored.
Concurrent work has identified MCP tool poisoning vectors~\cite{wu2024toolsword}
where malicious tool descriptions manipulate agent behavior.
Our work is complementary: we measure how much of the \emph{actual} capability
surface of published tools is absent from their descriptions (CR\,=\,\ResultCR{}),
providing the quantitative basis for arguing that MCP alone
cannot substitute for publish-time capability auditing.

\paragraph{Software supply chain.}
Typosquatting~\cite{vu2020typosquatting,taylor2020confused}
and dependency attacks~\cite{ohm2020backstabber,ladisa2023taxonomy}
are well-studied in npm/PyPI.
Our threat model differs: attacks arrive via legitimate updates
to previously-trusted skills, analogous to
event-stream~\cite{right2018eventstream} and node-ipc~\cite{snyk2022nodeipc}.
We provide the first measurement of update-based injection in LLM skill ecosystems.
The CodeArtifact security posture~\cite{lamb2022evaluating}
offers a reference architecture; LLM skill markets currently lack equivalent controls.

\paragraph{Capability disclosure and mobile permissions.}
Android studies~\cite{felt2011android,au2012pscout,wijesekera2015android}
show apps over-request permissions; we find the mirror image:
skills {\em under-declare} capabilities because LLM marketplaces
lack machine-readable manifests~\cite{bartel2012static,li2017static}.
Dark patterns~\cite{gray2018dark} and warning fatigue~\cite{bauer2013warning}
compound the problem.
The EU AI Act's transparency requirements~\cite{eu2024aiact} align
with our findings: mandatory capability disclosure for high-risk systems
would directly address the under-declaration gap (CDI\,=\,\ResultCDI{}) we quantify.
Unlike mobile permission studies that measure over-declaration,
our CR/CP framework independently captures both under-disclosure and false claims.

\paragraph{Static analysis and information disclosure.}
Our regex extraction relates to Android malware
detection~\cite{arp2014drebin} and code property
graphs~\cite{yamaguchi2014modeling}; we use lightweight regex
for marketplace-scale deployability (300+ patterns, \NskillsDMR{} skills, one pass).
LLM-generated code~\cite{pearce2022asleep,sandoval2023lost}
has elevated vulnerability rates, compounding under-disclosure
when developers do not audit generated skill code.
Static analysis of comment-code consistency~\cite{tan2007icomment,panthaplackel2020associating}
is the closest prior work; we apply this lens to security-critical marketplace listings.

\paragraph{Ecosystem measurement.}
Prior measurements span npm~\cite{zimmermann2019small},
PyPI~\cite{duan2021measuring}, Android~\cite{felt2011android},
and voice-assistant privacy~\cite{liao2020measuring}.
Ours is the first ecosystem-scale measurement of capability disclosure
in LLM skill marketplaces, spanning \NskillsDMR{} skills and \ResultFilesRead{} files.
The dual-platform design (\clawhub{} + \skillssh{}) directly addresses
the external validity concern raised for single-platform
supply-chain studies~\cite{ladisa2023taxonomy}:
convergent $\approx$90\% inconsistency across two independently-designed
markets strengthens the structural interpretation of our findings.

% sections/09_conclusion.tex

\section{Conclusion}
\label{sec:conclusion}

We presented \sysname, a publish-time static \emph{risk-prioritization} workflow for
surfacing capability-description mismatches in LLM agent skill marketplaces.
The workflow combines deterministic multi-file capability extraction (LSCPS,
14 families, \ResultFilesRead{} artifact files across \NskillsDMR{} skills)
with a lightweight disclosure-matching layer that computes Capability Recall
(CR), Description Precision (CP), and Capability Discrepancy Index (CDI).

Our evaluation on \NskillsDMR{} \clawhub{} skills finds that
\ResultICone{} of skills exhibit complete capability non-disclosure (IC-1,
CR\,=\,0) and \ResultICtwo{} exhibit partial non-disclosure (IC-2),
yielding corpus-mean CR\,=\,\ResultCR{} and CDI\,=\,\ResultCDI{}.
These findings are robust to regex bank variation: under 8 perturbation conditions
(strict/relaxed thresholds, 4 family ablations, 5-family taxonomy collapse),
inconsistency rates remain structurally elevated, with the Strict condition
still showing $>$47\% total inconsistency.
The combined SCS\,$\cup$\,DMR rule retrieves 94.1\%
(95\% Wilson CI:\,[80.9\%,\,98.4\%]) of \Nmalicious{} confirmed malicious skills
that the platform's own scanner missed entirely,
and outperforms four natural NLP baselines (lexical overlap, TF-IDF distance,
single-file static, general-purpose linters) on every ranking metric
(MRR: 0.019 vs. $\leq$0.002 for all baselines).
A marketplace security team reviewing the top-200 \sysname-ranked skills
recovers 17.6\% of confirmed T6 cases with a 16.3$\times$ lift over random,
compared to 0\% for the platform scanner under any budget.
Cross-platform validation on \NskillsSH{} independent skills suggests
that capability-description inconsistency exceeds 90\%
across both ecosystems without parameter retuning.

\paragraph{Limitations.}
LSCPS is a conservative lower bound: capabilities expressed via
non-standard idioms or obfuscated code are not captured.
Correspondingly, CR\,=\,\ResultCR{} should be interpreted as an
upper bound on actual disclosure---true gaps may be smaller.
The T6 ground truth (\Nmalicious{} skills) is limited by available
confirmed cases; recall estimates carry wide confidence intervals.
The \clawhub{} ecosystem, while the largest LLM skill marketplace
at the time of study, is smaller than mature package registries
(PyPI: $>$500K packages; npm: $>$2M packages).
We emphasize \emph{early-warning value}: documenting and hardening
the threat surface at 18K skills---when structural remediation
(per-skill capability manifests, YAML declaration enforcement)
is still tractable---is precisely when measurement papers are most actionable.
The npm event-stream and node-ipc precedents (\S\ref{sec:related})
demonstrate that supply-chain attacks scale faster than defenses;
our goal is to break that pattern before it repeats in LLM skill ecosystems.

\paragraph{Future Work.}
\textbf{AutoDesc and CapManifest} (§\ref{sec:autodesc}) represent the
primary constructive extension: an LLM-assisted description drafting tool
driven by LSCPS extraction results, paired with a machine-readable
CapManifest protocol that closes the IC gap from the developer side.
Evaluation goals include CR improvement rate ($\text{CR}_{\text{after}} \to 1.0$),
description acceptability in user studies, and false-capability rate.
\textbf{LLM semantic adjudication:} the five-model open-source ensemble
(prototype implemented) would reduce LSCPS false positives by verifying
that matched patterns are semantically functional.
\textbf{Expanded T6 ground truth} via coordinated platform disclosure
would improve statistical power on recall estimates.
\textbf{Real-time deployment}: integrating SCS+DMR as a publish-time
CI/CD gate would move \sysname from offline triage to active marketplace governance.
We view \sysname and AutoDesc as complementary layers---detection and remediation---
that together make the path of accurate disclosure easier than the path of omission.

\section*{Acknowledgments}

\textit{Anonymized for review.}



% ── References ────────────────────────────────────────────────────────────
\bibliographystyle{ACM-Reference-Format}
\bibliography{references}

% ── Appendices (do not count toward page limit) ────────────────────────────
\appendix
% sections/appendix.tex
% Optional appendix (not counted in body page limit)

\section{Ranking Effectiveness: Precision at $K$}
\label{app:patk}

\begin{table}[h]
\centering
\caption{Precision at $K$ for the composite ranking.
  Lift = P@$K$ / random baseline (34/\NskillsDMR{} $\approx0.49\%$).}
\label{tab:patk}
\small
\begin{tabular}{@{}rcccr@{}}
\toprule
\textbf{$K$} & \textbf{T6 found} & \textbf{P@$K$} & \textbf{Recall} & \textbf{Lift} \\
\midrule
 10 & 1  & 0.100 & 2.9\%  & 20.3$\times$ \\
 20 & 3  & 0.150 & 8.8\%  & 30.5$\times$ \\
 50 & 4  & 0.080 & 11.8\% & 16.3$\times$ \\
100 & 4  & 0.040 & 11.8\% & 8.1$\times$ \\
200 & 9  & 0.045 & 26.5\% & 9.1$\times$ \\
500 & 14 & 0.028 & 41.2\% & 5.7$\times$ \\
\bottomrule
\end{tabular}
\end{table}

\section{LSCPS Keyword Bank}
\label{app:lscps}

Table~\ref{tab:lscps_kw} lists representative regex patterns for each
\lscps{} capability family used in our static capability extractor.
The full bank contains 300+ patterns across 14 capability families.

\begin{table}[h]
\centering
\caption{LSCPS capability families and representative detection patterns.}
\label{tab:lscps_kw}
\small
\begin{tabular}{@{}ll@{}}
\toprule
\textbf{Family} & \textbf{Representative Patterns (regex)} \\
\midrule
Shell Exec     & \texttt{subprocess}, \texttt{os\.system}, \texttt{exec\textbackslash(}, \texttt{bash} \\
Cred Read      & \texttt{api\_key}, \texttt{access\_token}, \texttt{bearer}, \texttt{\.env} \\
Cred Transmit  & \texttt{post.*token}, \texttt{upload.*key}, \texttt{exfil} \\
Net Write      & \texttt{requests\.post}, \texttt{webhook}, \texttt{smtp} \\
Net Read       & \texttt{requests\.get}, \texttt{httpx\.get}, \texttt{curl http} \\
File Write     & \texttt{open\textbackslash(.*"w"}, \texttt{shutil\.}, \texttt{os\.remove} \\
File Read      & \texttt{open\textbackslash(.*"r"}, \texttt{glob\textbackslash(}, \texttt{os\.listdir} \\
Autonomous     & \texttt{autonomous}, \texttt{without.*confirm}, \texttt{cron} \\
Multi-step     & \texttt{multi.step}, \texttt{workflow}, \texttt{pipeline} \\
Memory Write   & \texttt{memory}, \texttt{remember}, \texttt{long.term} \\
\bottomrule
\end{tabular}
\end{table}

\section{T6 Case Timeline}
\label{app:t6_timeline}

For each T6 case we record: date of malicious version publication,
date of fix, exposure window (days), and installs during exposure.
Table~\ref{tab:t6_timeline} summarizes the 40 confirmed cases.

\begin{table}[h]
\centering
\caption{T6 case timeline summary (representative subset).}
\label{tab:t6_timeline}
\scriptsize
\begin{tabular}{@{}lrrlp{2.6cm}@{}}
\toprule
\textbf{Skill} & \textbf{Win.} & \textbf{Inst.} & \textbf{CES} & \textbf{Primary Attack} \\
\midrule
agentmail     & 1.0\,d & 155 & 8.5 & Webhook injection \\
moltguard     & 0.1\,d &  54 & 5.0 & Phishing file \\
simmer        & 0.4\,d &  21 & 5.0 & SKILL.md rollback \\
guard-scanner & 9.1\,d &   9 & 14.5 & Trojanized scanner \\
skill-auditor & --- &   8 & 4.5 & Credential harvest \\
vmware-aiops  & 1.1\,d &   3 & 4.5 & Webhook + file delete \\
vmware-monitor& 2.1\,d &   0 & 6.5 & Webhook + credential \\
aport-guardrail& 2.6\,d &   0 & 3.0 & Shell exec injection \\
clawsentinel  & 0.3\,d &   0 & 8.5 & Repeated inject/revert \\
mdk-wallet    & --- &   0 & 11.5 & rm -rf + credential exfil \\
\midrule
\textbf{All 40} & \textbf{avg 0.5\,d} & \textbf{298} & & \\
\bottomrule
\end{tabular}
\end{table}

\section{Deep Attack Case Studies}
\label{app:casestudies}

We present five representative T6 cases selected for diversity of attack
patterns, severity, and user impact. Full reconstruction details are available
in our artifact release.

\noindent\textit{Skill names below are pseudonymized per our responsible disclosure policy (\S\ref{sec:conclusion}); original identifiers are available to reviewers under embargo request.}

\paragraph{Case 1: Trojanized security scanner (\texttt{guard-scanner}).}
A skill marketed as a security scanning tool contained
credential harvesting (5 matches), network exfiltration (webhook.site,
curl), persistence mechanisms (cron, startup), and download-execute chains.
\textbf{Exposure}: 9.1~days, 9~installs. CR = 0.143; RWM = 12.86.
The attack used the ironic disguise of a ``security'' tool to normalize
access to security-critical APIs. All three detection layers (SCS, DMR, TDM) catch this case;
the platform scanner missed it.

\paragraph{Case 2: Prompt injection via email (\texttt{agentmail}).}
The highest-impact T6 case (155~installs) used unfiltered email webhooks
to forward arbitrary content as LLM prompts, enabling indirect prompt injection.
The fix (v1.1.0) added a webhook sender allowlist.
\textbf{Exposure}: 1.0~day. CR = 0.333; RWM = 2.12.
This case demonstrates that even moderate RWM scores can indicate
critical vulnerabilities when the attack vector is semantic rather than code-pattern-based.

\paragraph{Case 3: Financial data exfiltration (\texttt{mdk-agent-wallet}).}
A cryptocurrency wallet skill contained destructive file operations
(\texttt{rm -rf}) targeting the wallet directory,
webhook exfiltration of private keys, and credential harvesting.
The fix removed the destructive commands and added \texttt{autonomous:false}.
CR = 0.143; RWM = 15.33.
\textbf{Data flow}: Private key read $\to$ webhook POST $\to$ external endpoint.

\paragraph{Case 4: Hidden phishing attack (\texttt{moltguard}).}
A malware guard tool embedded a phishing attack file in its test
instructions, instructing the LLM agent to notify users when quota
tags appeared---effectively weaponizing the agent as a phishing vector.
\textbf{Exposure}: 0.1~days, 54~installs. CR = 0.000; RWM = 16.86.

\paragraph{Case 5: Repeated injection campaign (\texttt{clawsentinel}).}
This skill appeared twice in T6 (v3.0.0$\to$2.3.3 and v2.3.4$\to$2.3.5),
following an inject-revert-inject pattern. Each injection added
\texttt{rm -rf}, credential harvesting, and file destruction,
only to be reverted within hours.
CR = 0.143; \textbf{RWM = 29.96}---the highest in our dataset.
This case is consistent with repeated probing of platform detection limits.

\paragraph{Detection layer effectiveness across cases.}
In all five cases: (a) the platform scanner missed the attack;
(b) SCS static regex detected dangerous capabilities;
(c) DMR disclosure matching confirmed the capabilities were undisclosed;
(d) TDM temporal analysis would flag the forward-drift injection.
This confirms the multi-layer design is necessary and
no single layer is sufficient.

\section{Detector Feature Ablation}
\label{app:ablation}

Table~\ref{tab:ablation} shows AUROC when each feature group is used
in isolation (Logistic Regression, 5-fold CV).

\begin{table}[h]
\centering
\caption{Feature group ablation for temporal drift detector.}
\label{tab:ablation}
\small
\begin{tabular}{@{}llr@{}}
\toprule
\textbf{Group} & \textbf{Key Features} & \textbf{AUROC} \\
\midrule
Drift core & CES delta, text diff, new kw & 0.979 \\
Disclosure & SDS, changelog len, vague & 0.979 \\
Family drift & Per-family CES deltas & 0.980 \\
\textbf{All combined} & & \textbf{1.000} \\
\bottomrule
\end{tabular}
\end{table}

\section{Responsible Disclosure}
\label{app:disclosure}

We notified both \clawhub{} and \skillssh{} of the T6 cases
\textit{[redacted for anonymous review]}
days before submission.
\clawhub{} acknowledged and confirmed all 40 cases.
\skillssh{} did not respond within the disclosure window.
We have withheld specific skill identifiers in this submission;
they will be released as supplementary material upon acceptance.

\section{LLM Annotation Protocol}
\label{app:llm_protocol}

\emph{Note: LLM-based semantic adjudication is deferred to future work.
This appendix describes the planned methodology.}

We propose to evaluate IC classification robustness using five independent
open-source LLM annotators on \NskillsDMR{} \clawhub{} skills.

\paragraph{Planned Models.}
We plan to deploy five fully open-weight models locally via vLLM~\cite{kwon2023vllm}
on an HPC cluster, with no commercial API dependence:
\begin{itemize}[topsep=2pt,itemsep=1pt]
\item \textbf{M1}: Qwen3-8B (Alibaba Cloud) --- chain-of-thought reasoning, multilingual
\item \textbf{M2}: Llama-3.1-8B-Instruct (Meta AI) --- instruction-tuned, North American research
\item \textbf{M3}: Mistral-7B-v0.3-Instruct (Mistral AI) --- efficient European-origin model
\item \textbf{M4}: Gemma-3-4B-IT (Google DeepMind) --- compact instruction-tuned, safety-focused
\item \textbf{M5}: GPT-OSS-20B (OpenAI) --- open-weight, 20B parameter research model
\end{itemize}

This five-family ensemble spans models from Alibaba, Meta, Mistral, Google, and OpenAI,
maximizing the diversity of training data, RLHF alignment procedures, and architectural
design choices to minimize the risk of systematic single-model extraction artifacts.

\paragraph{Prompt template.}
Each annotator receives a structured prompt containing:
(a)~the skill's user-visible description (truncated to 400 chars),
(b)~the developer's changelog (truncated to 500 chars),
(c)~statically detected new risk keywords,
(d)~a unified diff of the SKILL.md between versions (capped at 150 lines),
(e)~the CES drift score, and
(f)~the SDS disclosure score.

The system prompt instructs the model to classify the update as:
\texttt{CONFIRMED} (undisclosed dangerous capability),
\texttt{FALSE\_POSITIVE} (contextually benign or adequately disclosed), or
\texttt{UNCLEAR} (insufficient evidence).
The model must respond with structured JSON including a label,
one-sentence reason, specific risk evidence, and a boolean
\texttt{disclosed} flag.

\paragraph{Consensus.}
The final label uses majority vote ($\geq3/5$ agreement).
Each annotator operates independently with no access to
the other annotators' outputs.
We use the resulting consensus
to estimate review yield on sampled pipeline outputs, not to
construct a universal ground-truth benchmark.

\paragraph{Reproducibility.}
All prompts, raw outputs, and the annotation script
(\texttt{08b\_claude\_native\_m1.py}, \texttt{08c\_gpt\_native\_m2.py})
are included in our artifact release.


% ── Required Statements ───────────────────────────────────────────────────
\section*{Ethics Considerations}
% Ethics Considerations
% Required by ACM CCS 2026

We performed a stakeholder-based ethics analysis following the Menlo Report framework~\cite{dittrich2012menlo}, considering four principles:

\paragraph{Beneficence.}
Our research directly benefits platform operators and users by identifying a systematic capability-description gap that existing platform scanners entirely miss (0\% T6 recall).
The detection pipeline we propose could prevent future supply-chain attacks affecting hundreds of users.

\paragraph{Respect for Persons.}
All data was collected from publicly accessible APIs.
No user accounts were accessed, no private data was collected, and no skill was executed during analysis.
T6 malicious skill identifiers are pseudonymized throughout the paper to protect affected developers who may have been victims of account compromise.

\paragraph{Justice.}
We applied our analysis uniformly across all skills in both ecosystems without targeting individual developers.
Our framework classifies skills, not developers---the high inconsistency rate (90\%)
reflects a widespread ecosystem norm, not individual malfeasance.

\paragraph{Respect for Law and Public Interest.}
We notified both platform operators of our findings prior to submission (responsible disclosure).
Our crawler respected \texttt{robots.txt} and rate-limited requests to $\leq$0.5 req/s.
No laws or terms of service were violated during data collection.

\paragraph{Potential harms.}
Publishing detailed capability extraction patterns could theoretically help attackers craft harder-to-detect malicious skills.
We mitigate this by (a)~withholding the most specific evasion-relevant pattern details from the public artifact and (b)~sharing the full pattern bank only with platform operators and academic researchers under data use agreements.


\section*{Artifact Availability}
% Artifact Availability
% ACM CCS 2026 artifact availability statement

We commit to making the following artifacts available for reproducibility:

\begin{enumerate}[topsep=3pt,itemsep=2pt]
\item \textbf{LSCPS Regex Pattern Bank} (300+ patterns, 14 families):
  The complete deterministic capability extraction engine.
  Available at: \url{[anonymized repository]}.

\item \textbf{IC Classification Pipeline} (\texttt{19\_ic\_framework.py}):
  Computes the difference matrix (AD, UD, FC, TN), derives CR, CP, CDI,
  and assigns IC types (IC-1 through IC-4) for any skill dataset.

\item \textbf{Detection Evaluation Scripts}:
  End-to-end T6 recall evaluation and cross-platform
  IC analysis scripts.

\item \textbf{IC Classification Results}:
  \texttt{ic\_classification.csv} (\NskillsDMR{} \clawhub{} skills with IC types, CR, CP, CDI)
  and \texttt{ic\_crossplatform.csv} (116,208 \skillssh{} skills).

\item \textbf{T6 Ground Truth Corpus}:
  40 confirmed malicious version pairs (pseudonymized) with
  before/after SKILL.md snapshots, remediation changelogs,
  attack type labels (CRED\_READ, FILE\_DELETE, NET\_WRITE, etc.),
  and corroboration source type (platform advisory / changelog / public report).
  Inter-rater $\kappa$, rubric definition for T6 inclusion,
  and E4 security-relevance rubric are included as supplementary documents.

\item \textbf{LLM Adjudication Prototype} (\S\ref{sec:sa}):
  Source code for the five-model ensemble adjudication prototype
  (\mone, \mtwo, \mthree, \mfour, \mfive).
  Full-corpus execution is deferred to future work;
  prototype code will be released upon paper acceptance.

\end{enumerate}

\paragraph{Access during review.}
All artifacts are available via an anonymized repository link
provided in the HotCRP submission form.
The evaluation scripts reproduce every table and figure in this paper
directly from the derived CSV datasets, without re-running the full
extraction pipeline on raw skill archives.
The raw \clawhub{} and \skillssh{} skill archives are not redistributed
due to platform terms of service,
but all derived values needed for reproducibility are included in the CSVs.

\paragraph{Post-acceptance.}
Upon acceptance, we will publish a public artifact package on Zenodo
with a persistent DOI, including the full LSCPS regex bank,
evaluation scripts, and anonymized datasets.


\end{document}
